% !TEX program = xelatex
% Biên dịch bằng XeLaTeX (khuyến nghị) để hỗ trợ tiếng Việt.
% Trên Overleaf: Menu -> Settings -> Compiler -> XeLaTeX.
\documentclass[a4paper,13pt,twoside]{report}
\usepackage{scrextend}
\changefontsizes{13pt}

\usepackage{fontspec}        % cần XeLaTeX/LuaLaTeX
\usepackage{polyglossia}
\setmainlanguage{vietnamese}
\setotherlanguage{english}
% Font kiểu Times New Roman, hỗ trợ tiếng Việt — nạp theo tên file để chạy được
% cả trên máy local (TeX Live) lẫn Overleaf, không phụ thuộc font hệ thống.
\setmainfont{texgyretermes}[
  Extension      = .otf,
  UprightFont    = *-regular,
  BoldFont       = *-bold,
  ItalicFont     = *-italic,
  BoldItalicFont = *-bolditalic,
]

\usepackage[top=2cm, bottom=2cm, left=3.5cm, right=2.5cm]{geometry}
\usepackage{indentfirst}     % thụt đầu dòng cả đoạn đầu tiên
\usepackage{setspace}        % giãn dòng
\usepackage{graphicx}
\usepackage{float}
\usepackage[table]{xcolor}
\usepackage{amsmath,amssymb}
\usepackage{hyperref}
\hypersetup{colorlinks=true, linkcolor=black, urlcolor=black, citecolor=black}
\usepackage{xurl}            % cho phép ngắt dòng URL/đường dẫn tại mọi vị trí
\usepackage{booktabs}
\usepackage{longtable}
\usepackage{array}
\usepackage{tabularx}
\usepackage{multirow}
\usepackage{makecell}
\usepackage{enumitem}
\usepackage{fancyhdr}
\usepackage{titlesec}
\usepackage{titletoc}
\usepackage[font=small,labelfont=bf]{caption}
\usepackage[most]{tcolorbox}
\usepackage{tikz}
\usetikzlibrary{shapes.geometric, arrows.meta, positioning, fit, backgrounds, calc}

% ----- Màu dùng cho các nhóm nghiên cứu và hộp nhận xét -----
\definecolor{gSafe}{HTML}{C0392B}
\definecolor{gOrch}{HTML}{2471A3}
\definecolor{gComm}{HTML}{1E8449}
\definecolor{gApp}{HTML}{B9770E}
\definecolor{boxbg}{HTML}{F4F6F7}

\newtcolorbox{nhanxet}[1][Nhận xét]{breakable,colback=boxbg,colframe=black!45,
  boxrule=0.5pt,arc=2pt,left=6pt,right=6pt,top=4pt,bottom=4pt,
  fonttitle=\bfseries\small,title=#1}

% ----- Kiểu cột bảng -----
\newcolumntype{L}[1]{>{\raggedright\arraybackslash}p{#1}}
\newcolumntype{C}[1]{>{\centering\arraybackslash}p{#1}}
\newcolumntype{Y}{>{\raggedright\arraybackslash}X}

% ----- Viết tắt -----
\newcommand{\MAS}{MAS}
\newcommand{\LLM}{LLM}

% ----- Giãn dòng, thụt đầu dòng (theo chuẩn đồ án) -----
\onehalfspacing
\setlength{\parskip}{6pt}
\setlength{\parindent}{15pt}
% Không kéo giãn khoảng cách giữa các khối chỉ để lấp đầy chiều cao trang.
\raggedbottom
% Tiếng Việt không ngắt-vần được: cho phép giãn khoảng trắng để tránh tràn lề
\setlength{\emergencystretch}{3em}
\setcounter{secnumdepth}{3}     % đánh số tới subsubsection
\counterwithin{figure}{chapter} % Hình 1.1, 1.2, ...
\counterwithin{table}{chapter}  % Bảng 1.1, 1.2, ...
\setlist{nosep}

% ----- Cho phép bảng/hình trôi linh hoạt để tránh khoảng trắng lớn -----
\renewcommand{\topfraction}{0.9}
\renewcommand{\bottomfraction}{0.8}
\renewcommand{\textfraction}{0.07}
\renewcommand{\floatpagefraction}{0.75}
\setcounter{topnumber}{3}
\setcounter{bottomnumber}{2}
\setcounter{totalnumber}{5}
% Trang chỉ có bảng/hình: canh lên đầu trang thay vì canh giữa
\makeatletter
\setlength{\@fptop}{0pt}
\setlength{\@fpsep}{12pt plus 2pt}
\setlength{\@fpbot}{0pt plus 1fil}
\makeatother

% ----- Header/Footer -----
\setlength{\headheight}{15pt}
\pagestyle{fancy}
\fancyhf{}
\fancyhead[RE,LO]{\leftmark}
\fancyfoot[RE,LO]{\thepage}
\renewcommand{\chaptermark}[1]{\markboth{#1}{}}

% ----- Định dạng tiêu đề chương/mục theo chuẩn HUST -----
\titleformat{\chapter}[hang]{\centering\bfseries}{CHƯƠNG \thechapter.\ }{0pt}{}
\titlespacing*{\chapter}{0pt}{-20pt}{20pt}

\titleformat{\section}[hang]{\bfseries}{\thechapter.\arabic{section}\quad}{0pt}{}
\titlespacing{\section}{0pt}{\parskip}{0.5\parskip}

\titleformat{\subsection}[hang]{\bfseries}{\thechapter.\arabic{section}.\arabic{subsection}\quad}{0pt}{}
\titlespacing{\subsection}{15pt}{\parskip}{0.5\parskip}

\renewcommand\thesubsubsection{\alph{subsubsection}}
\titleformat{\subsubsection}[hang]{\bfseries}{\alph{subsubsection})\ }{0pt}{}
\titlespacing{\subsubsection}{30pt}{\parskip}{0.5\parskip}

% polyglossia ghi đè tên mục khi vào tài liệu, nên đặt lại trong captionsvietnamese
\gappto\captionsvietnamese{\renewcommand*\contentsname{MỤC LỤC}\renewcommand*\bibname{TÀI LIỆU THAM KHẢO}}

% ----- Định dạng mục lục (dấu chấm dẫn, tiền tố CHƯƠNG) -----
\titlecontents{chapter}[0em]
  {\addvspace{0.3cm}\bfseries}
  {CHƯƠNG \thecontentslabel.\quad}
  {}
  {\titlerule*[0.3pc]{.}\contentspage}
\titlecontents{section}[1.4em]
  {\addvspace{0.1cm}}
  {\thecontentslabel\quad}
  {}
  {\titlerule*[0.3pc]{.}\contentspage}
\titlecontents{subsection}[2.8em]
  {}
  {\thecontentslabel\quad}
  {}
  {\titlerule*[0.3pc]{.}\contentspage}

\begin{document}

% ===================== TRANG BÌA =====================
\begin{titlepage}
\thispagestyle{empty}
\begin{center}

{\textbf{\large{ĐẠI HỌC BÁCH KHOA HÀ NỘI}}}\\[4cm]

{\textbf{\huge{BÁO CÁO BÀI TẬP LỚN}}}\\[1cm]
{\textbf{\Large{Multi-Agent Systems và Collaborative Reasoning}}}\\[1cm]

{\textbf{\large{TRẦN XUÂN HƯƠNG}}}\\
{\textbf{\large{VŨ TÙNG DƯƠNG}}}\\
{\textbf{\large{PHẠM VIỆT DŨNG}}}\\
{\textbf{\large{LÊ VĂN ĐỨC}}}\\[0.5cm]

{\textbf{\large{Học phần: Xử lý ngôn ngữ tự nhiên}}}\\

\vspace{2cm}
\begin{table}[H]
\centering
\resizebox{\textwidth}{!}{%
\begin{tabular}{ll}
\multicolumn{1}{c}{\textbf{Giảng viên hướng dẫn:}} & PGS.TS. Lê Thanh Hương \hspace{0.5cm} \underline{\hspace{1.8cm}}  \\[0.5cm]
  & \multicolumn{1}{r}{Chữ kí GVHD}     \\[0.5cm]
\textbf{Mã học phần:} & IT4772                 \\[0.5cm]
\textbf{Trường:} & Công nghệ Thông tin và Truyền thông \\[3cm]
\multicolumn{2}{c}{\textbf{HÀ NỘI, 10/2026}}
\end{tabular}%
}
\end{table}
\end{center}
\end{titlepage}

\pagenumbering{roman}

% ===================== LỜI CẢM ƠN =====================
\thispagestyle{fancy}
\begin{center}
\LARGE\textbf{LỜI CẢM ƠN}
\end{center}
\vspace{1cm}
Nhóm chúng em xin gửi lời cảm ơn chân thành tới \textbf{PGS.TS. Lê Thanh Hương} đã tận tình giảng
dạy, góp ý và định hướng cho nhóm trong suốt quá trình học tập và thực hiện bài tập lớn của học phần
Xử lý ngôn ngữ tự nhiên. Chúng em cũng xin cảm ơn các thầy cô trong Trường Công nghệ Thông tin và
Truyền thông -- Đại học Bách khoa Hà Nội đã trang bị cho chúng em những kiến thức nền tảng quý báu.
Do thời gian và kiến thức còn hạn chế, báo cáo không tránh khỏi thiếu sót; nhóm rất mong nhận được
sự góp ý của cô để hoàn thiện hơn. Chúng em xin chân thành cảm ơn!
\newpage

% ===================== TÓM TẮT =====================
\thispagestyle{fancy}
\begin{center}
\Large\textbf{TÓM TẮT NỘI DUNG BÁO CÁO}
\end{center}
\vspace{1cm}
LLM-based Multi-Agent Systems (\MAS) cho phép nhiều agent chuyên biệt cùng phân rã, trao đổi và
tổng hợp lập luận để giải các bài toán mà một mô hình đơn lẻ khó xử lý tốt. Báo cáo khảo sát, phân
tích và so sánh tám nghiên cứu mới nhất trên arXiv (tháng 5 -- tháng 9/2026) về chủ đề
\emph{Multi-Agent Systems và Collaborative Reasoning}, phân loại thành bốn hướng:
(i) safety \& robustness (MIRAGE, TIDE); (ii) orchestration \& control (ProgRouter, HalluProp);
(iii) communication \& interaction (DAIN, AGENT-RADAR); (iv) domain applications (MAC-RRG,
UrbanAgent). Với mỗi nghiên cứu, báo cáo trình bày bài toán, phương pháp, thiết lập thực nghiệm,
kết quả định lượng chính và đánh giá phản biện; sau đó so sánh có hệ thống theo phương pháp, dữ
liệu, kết quả, chi phí, ưu điểm, hạn chế và phạm vi áp dụng.

Ba nhận định chính: (1) \emph{communication graph} là đối tượng trung tâm của hầu hết các nghiên
cứu; (2) sự lan truyền của lỗi, hallucination và thông tin sai qua các communication edge là rủi ro
cốt lõi, và thảo luận nhiều vòng có xu hướng khuếch đại sai lệch thay vì sửa sai; (3) khái niệm
``agent'' đang được dùng rất rộng, gây khó khăn cho việc so sánh. Cuối cùng, báo cáo chỉ ra các
khoảng trống nghiên cứu và đề xuất một số hướng phát triển.

\vspace{6pt}
\noindent\textbf{Từ khóa:} multi-agent systems, collaborative reasoning, large language models,
agent orchestration, communication graph, AI safety, hallucination propagation.
\newpage

% ===================== MỤC LỤC, DANH MỤC =====================
\tableofcontents
\newpage

\renewcommand{\listfigurename}{DANH MỤC HÌNH VẼ}
{\let\oldnumberline\numberline
\renewcommand{\numberline}{Hình~\oldnumberline}
\listoffigures}
\newpage

\renewcommand{\listtablename}{DANH MỤC BẢNG BIỂU}
{\let\oldnumberline\numberline
\renewcommand{\numberline}{Bảng~\oldnumberline}
\listoftables}
\newpage

% ===================== DANH MỤC TỪ VIẾT TẮT =====================
\thispagestyle{fancy}
\begin{center}
\Large\textbf{DANH MỤC THUẬT NGỮ VÀ TỪ VIẾT TẮT}
\end{center}
\vspace{0.3cm}
\renewcommand{\arraystretch}{1.3}
\begin{center}
\begin{tabular}{p{2.0cm} >{\raggedright\arraybackslash}p{4.4cm} >{\raggedright\arraybackslash}p{7.0cm}}
\toprule
\textbf{Viết tắt} & \textbf{Tiếng Anh} & \textbf{Ý nghĩa} \\
\midrule
LLM   & Large Language Model           & Mô hình ngôn ngữ lớn \\
MAS   & Multi-Agent System             & Hệ thống gồm nhiều agent cùng phối hợp giải bài toán \\
agent & Agent                          & Thực thể dùng LLM (hoặc mô-đun) nhận đầu vào, suy luận và tạo đầu ra/hành động \\
DAG   & Directed Acyclic Graph         & Đồ thị có hướng không chu trình \\
RAG   & Retrieval-Augmented Generation & Sinh văn bản có bổ sung tri thức truy xuất \\
KG    & Knowledge Graph                & Đồ thị tri thức \\
MoE   & Mixture-of-Experts             & Kiến trúc chọn một tập con expert cho mỗi đầu vào \\
PID   & Partial Information Decomposition & Phân rã thông tin thành uniqueness, redundancy, synergy \\
SFT   & Supervised Fine-Tuning         & Tinh chỉnh có giám sát \\
GRPO  & Group Relative Policy Optimization & Thuật toán học tăng cường dùng advantage chuẩn hoá theo nhóm \\
GNN   & Graph Neural Network           & Mạng nơ-ron trên đồ thị \\
SPA   & Selective Prompt Anchoring     & Kỹ thuật attention steering nhấn mạnh một phần prompt \\
CIA   & Communication Inference Attack & Tấn công suy ra topology của MAS \\
AUROC & Area Under the ROC Curve       & Diện tích dưới đường cong ROC \\
POI   & Point of Interest              & Điểm quan tâm trên bản đồ \\
\bottomrule
\end{tabular}
\end{center}
\renewcommand{\arraystretch}{1}
\newpage

\pagenumbering{arabic}

\chapter{GIỚI THIỆU}
\label{chap:intro}

\section{Bối cảnh}

Các mô hình ngôn ngữ lớn (\LLM) đã đạt năng lực đáng kể trong hiểu và sinh ngôn ngữ, lập luận toán học, lập trình và trả lời câu hỏi. Tuy nhiên, một mô hình đơn lẻ vẫn gặp nhiều giới hạn: dễ sinh hallucination, khó tự kiểm tra lại lập luận, khó xử lý các tác vụ dài nhiều bước và khó kết hợp nhiều nguồn tri thức hoặc nhiều phương thức dữ liệu khác nhau. Một hướng tiếp cận tự nhiên là tổ chức \emph{nhiều} agent dựa trên \LLM{} thành một hệ thống, trong đó mỗi agent đảm nhận một vai trò (người lập kế hoạch, người lập trình, người phản biện, chuyên gia miền,~\ldots) và các agent trao đổi thông điệp để cùng đi đến lời giải. Đây là ý tưởng cốt lõi của LLM-based Multi-Agent Systems (\MAS) và của \emph{collaborative reasoning}.

Từ các khung như CAMEL~\cite{camel}, AutoGen~\cite{autogen}, MetaGPT~\cite{metagpt} đến cơ chế multi-agent debate~\cite{du2024debate}, cộng đồng nghiên cứu đã chỉ ra rằng sự cộng tác có thể cải thiện tính đúng đắn sự kiện và khả năng lập luận. Các công trình tiếp theo như GPTSwarm~\cite{gptswarm}, G-Designer~\cite{gdesigner} hay AgentPrune~\cite{agentprune} coi \MAS{} như một \emph{communication graph} có thể tối ưu được. Đến năm 2026, khi \MAS{} được triển khai rộng rãi hơn, các câu hỏi nghiên cứu trở nên thực tế và khó hơn:
\begin{itemize}
  \item \textbf{Hiệu quả:} chạy nhiều agent tốn kém (token, năng lượng GPU, độ trễ). Làm sao chọn đúng mô hình cho đúng agent ở đúng thời điểm?
  \item \textbf{Luồng thông tin:} lịch sử hội thoại giữa các agent tăng rất nhanh. Làm sao để mỗi agent chú ý đúng vào thông tin liên quan?
  \item \textbf{An toàn và độ tin cậy:} một agent sai, bị thao túng hoặc nói dối có thể kéo cả hệ thống đi sai. Cấu trúc hệ thống có thể bị kẻ tấn công suy ra và khai thác.
  \item \textbf{Ứng dụng:} làm sao áp dụng mô hình cộng tác vào các miền chuyên sâu như y tế hay đô thị, nơi dữ liệu đa phương thức và mâu thuẫn giữa các nguồn là phổ biến?
\end{itemize}

\section{Mục tiêu và phạm vi khảo sát}

Báo cáo thực hiện theo hình thức \emph{tiểu luận} của học phần IT4772, với chủ đề số 3: \emph{Multi-Agent Systems và Collaborative Reasoning}. Mục tiêu gồm:
\begin{enumerate}
  \item Tổng quan và phân loại các nghiên cứu theo hướng tiếp cận chính.
  \item Trình bày đóng góp của từng nghiên cứu tiêu biểu.
  \item So sánh, đánh giá có hệ thống về phương pháp, dữ liệu, kết quả, ưu điểm, hạn chế và phạm vi áp dụng.
  \item Chỉ ra khoảng trống nghiên cứu và đề xuất hướng phát triển.
\end{enumerate}

\paragraph{Phương pháp chọn bài.} Nhóm truy vấn kho arXiv với các bài có cụm từ ``multi-agent'' và ``collaborative reasoning'' trong phần tóm tắt, sắp xếp theo ngày nộp mới nhất, và chọn tám bài mới nhất (nộp từ 28/05/2026 đến 29/09/2026). Cách chọn này ưu tiên tính \emph{thời sự}, phản ánh các hướng mà cộng đồng đang quan tâm tại thời điểm thực hiện báo cáo. Hệ quả là tập bài khá đa dạng: có bài phương pháp tổng quát, có bài về an toàn, và có bài ứng dụng chuyên ngành. Sự đa dạng này vừa là thách thức cho việc so sánh trực tiếp, vừa cho phép nhìn được bức tranh rộng của lĩnh vực. Danh sách tám bài được liệt kê ở Bảng~\ref{tab:paperlist}.

\begin{table}[!htbp]
\centering
\small
\caption{Tám nghiên cứu được khảo sát (sắp theo ngày nộp lên arXiv, mới nhất trước).}
\label{tab:paperlist}
\begin{tabularx}{\textwidth}{c L{2.75cm} Y c}
\toprule
\textbf{\#} & \textbf{Tên gọi} & \textbf{Tiêu đề} & \textbf{arXiv} \\
\midrule
1 & MIRAGE~\cite{p1mirage} & Concealing LLM-Based Multi-Agent Topology via Phantom Structure Injection & 2609.37567 \\
2 & MAC-RRG~\cite{p2macrrg} & MAC-RRG: Iterative Multi-Agent Collaboration for X-ray Radiology Report Generation & 2609.26124 \\
3 & ProgRouter~\cite{p3progrouter} & ProgRouter: Online Progress-Guided Orchestration for Multi-Agent LLM Workflows under Quality-Cost Tradeoffs & 2608.25992 \\
4 & TIDE~\cite{p4tide} & When Truth Is Distributed: Misinformation Derails Collective Fact Recovery in LLM-Based Multi-Agent Systems & 2608.03421 \\
5 & HalluProp~\cite{p5halluprop} & Before Agents Speak: Pre-hoc Failure Risk Inference in Multi-Agent Systems & 2607.26836 \\
6 & UrbanAgent~\cite{p6urban} & Multi-Agent Collaborative Reasoning with Tool-Augmented Evidence for Urban Region Profiling & 2607.13558 \\
7 & DAIN~\cite{p7dain} & DAIN: Dynamic Agent-Based Interaction Network for Efficient and Collaborative Multimodal Reasoning & 2606.30189 \\
8 & AGENT-RADAR~\cite{p8radar} & Enhancing Multi-Agent Communication through Attention Steering with Context Relevance & 2605.30136 \\
\bottomrule
\end{tabularx}
\end{table}

\paragraph{Lưu ý về nguồn.} Phần lớn các bài là bản tiền ấn phẩm (preprint) chưa qua bình duyệt; chỉ UrbanAgent ghi rõ được chấp nhận tại KDD~2026. Một số bảng số liệu và hình trong bản PDF bị biến dạng khi trích xuất văn bản; những số liệu mà nhóm không chắc chắn đều được ghi chú rõ hoặc không đưa vào báo cáo.

\section{Đóng góp của báo cáo}

\begin{itemize}
  \item Đề xuất một \textbf{khung phân loại} bốn nhóm cho tám nghiên cứu, kèm \emph{ma trận multi-label} cho thấy các bài thường thuộc nhiều nhóm cùng lúc, và một \textbf{lược đồ vòng đời} \MAS{} cho biết mỗi phương pháp can thiệp vào giai đoạn nào (trước khi chạy, trong khi chạy, khi giao tiếp, khi phơi bày kết quả).
  \item Tóm tắt có cấu trúc và \textbf{đánh giá phản biện} từng nghiên cứu, trong đó chỉ ra cả những điểm yếu mà tác giả không nêu (ví dụ: không báo cáo phương sai, số liệu không nhất quán giữa các bảng, điều chỉnh siêu tham số trên tập kiểm thử).
  \item \textbf{So sánh hệ thống} theo bảy chiều và rút ra các xu hướng xuyên suốt.
  \item Xác định \textbf{khoảng trống nghiên cứu} và đề xuất hướng phát triển cụ thể.
\end{itemize}

\section{Cấu trúc báo cáo}

Chương~\ref{chap:background} trình bày kiến thức nền về \MAS{} dựa trên \LLM{} và collaborative reasoning. Chương~\ref{chap:taxonomy} đề xuất khung phân loại. Chương~\ref{chap:papers} trình bày chi tiết tám nghiên cứu theo bốn nhóm. Chương~\ref{chap:comparison} so sánh và đánh giá có hệ thống. Chương~\ref{chap:gaps} thảo luận khoảng trống nghiên cứu và hướng phát triển. Chương~\ref{chap:conclusion} kết luận. Chương~\ref{chap:contrib} trình bày đóng góp của từng thành viên.

\chapter{KIẾN THỨC NỀN}
\label{chap:background}

\section{LLM-based agent}

Một \emph{agent} dựa trên \LLM{} là một thực thể nhận đầu vào (câu hỏi, quan sát, thông điệp từ agent khác), suy luận bằng \LLM{} và tạo ra hành động hoặc thông điệp đầu ra. Agent thường được định nghĩa bởi ba thành phần: (i) \emph{backbone LLM}; (ii) \emph{role/system prompt} mô tả chuyên môn và trách nhiệm; (iii) tùy chọn \emph{công cụ} (tool) như tìm kiếm web, trình thông dịch mã, retrieval tài liệu. Mô hình ReAct~\cite{react} là khuôn mẫu phổ biến, trong đó agent xen kẽ giữa bước \emph{reasoning} và bước \emph{action} (gọi công cụ), rồi cập nhật dựa trên \emph{quan sát} trả về.

\section{MAS như một communication graph}

Phần lớn các nghiên cứu gần đây mô hình hóa \MAS{} như một đồ thị có hướng
\begin{equation}
  \mathcal{G} = (\mathcal{V}, \mathcal{E}),
\end{equation}
trong đó mỗi đỉnh $v_i \in \mathcal{V}$ là một agent và cạnh $(v_j \to v_i) \in \mathcal{E}$ cho biết đầu ra của $v_j$ được chuyển làm ngữ cảnh đầu vào cho $v_i$. Với câu hỏi $q$, role prompt $p_i$ và tập đầu ra của các agent đứng trước $O_i = \{r_j : (v_j \to v_i)\in\mathcal{E}\}$, agent $v_i$ sinh ra
\begin{equation}
  r_i = \mathrm{LLM}(p_i,\; q,\; O_i).
  \label{eq:agent}
\end{equation}
Hệ có thể chạy nhiều \emph{vòng} (round): ở vòng $t$, agent có thể đọc toàn bộ đầu ra của các vòng trước. Một decision agent (decision node) tổng hợp và trả lời. Cách nhìn này thống nhất được nhiều kiểu tổ chức khác nhau và là nền tảng của các phương pháp \emph{topology optimization} như GPTSwarm~\cite{gptswarm}, G-Designer~\cite{gdesigner} hay \emph{cắt tỉa} giao tiếp như AgentPrune~\cite{agentprune}.

\section{Các mô hình cộng tác phổ biến}

Bảng~\ref{tab:paradigms} tóm tắt các mô hình cộng tác thường gặp; các nghiên cứu được khảo sát sử dụng hầu hết các mô hình này.

\begin{table}[!htbp]
\centering
\small
\caption{Các mô hình tổ chức cộng tác trong \MAS{} dựa trên \LLM.}
\label{tab:paradigms}
\begin{tabularx}{\textwidth}{L{3.1cm} Y L{3.2cm}}
\toprule
\textbf{Mô hình} & \textbf{Mô tả} & \textbf{Ví dụ} \\
\midrule
Debate & Nhiều agent đưa ra câu trả lời, đọc câu trả lời của nhau và cập nhật qua nhiều vòng; kết quả là đồng thuận hoặc bỏ phiếu đa số. & Du et al.~\cite{du2024debate}; TIDE (\#4) \\
Phân cấp / coordinator & Một coordinator agent (coordinator/orchestrator) lập kế hoạch, giao việc cho các agent thực thi và theo dõi tiến độ. & Magentic-One~\cite{magenticone}; ProgRouter (\#3) \\
Dây chuyền vai trò & Các agent chuyên biệt xử lý tuần tự như một quy trình phần mềm (phân tích -- thiết kế -- lập trình -- kiểm thử). & MetaGPT~\cite{metagpt}; MAC-RRG (\#2) \\
Đồ thị tùy ý / tối ưu được & Topology là đồ thị (ngẫu nhiên, phân lớp, đầy đủ) có thể học hoặc cắt tỉa. & GPTSwarm, G-Designer, AgentPrune; HalluProp (\#5), AGENT-RADAR (\#8) \\
Router & Một bộ điều khiển chọn agent hoặc mô hình nào được kích hoạt cho từng đầu vào/bước. & MoE; ProgRouter (\#3), DAIN (\#7) \\
\bottomrule
\end{tabularx}
\end{table}

\section{Các thách thức chính của MAS}

\subsection{Chi phí và hiệu quả}
Mỗi agent là một hoặc nhiều lần gọi \LLM, và ngữ cảnh tích lũy qua các vòng làm số token tăng nhanh. Chi phí có thể đo bằng token, tiền API, năng lượng GPU hay độ trễ. Các phương pháp \emph{model routing} như RouteLLM hay các phương pháp \emph{cascade} truyền thống chỉ ra quyết định một lần cho mỗi câu hỏi, trong khi quy trình multi-agent cần chuỗi quyết định phụ thuộc trạng thái~\cite{p3progrouter}.

\subsection{Quá tải ngữ cảnh}
Khi lịch sử hội thoại dài, \LLM{} có xu hướng bỏ sót thông tin nằm ở giữa ngữ cảnh -- hiện tượng ``lost in the middle''~\cite{lostmiddle}. Trong \MAS{}, các agent còn lặp lại nội dung của nhau, khiến mô tả nhiệm vụ gốc và bằng chứng quan trọng bị ``chôn vùi''. Các giải pháp thường dùng là tóm tắt định kỳ (compaction) hoặc cắt tỉa agent/cạnh (AgentPrune, AgentDropout), nhưng cả hai đều làm mất thông tin~\cite{p8radar}.

\subsection{Lan truyền lỗi, hallucination và thông tin sai}
Theo phương trình~\eqref{eq:agent}, đầu ra của một agent trở thành đầu vào của agent khác. Do đó, một hallucination hoặc một thông tin sai cố ý có thể \emph{lan truyền} theo các cạnh, được các agent khác tiếp nhận, lặp lại và khuếch đại, dẫn tới cascading failure ở cấp hệ thống~\cite{p5halluprop,p4tide}. Nghiên cứu tâm lý học nhóm ở người về ``hidden profile''~\cite{stasser1985} cho thấy các nhóm có xu hướng ưu tiên thông tin được chia sẻ chung và bỏ qua thông tin riêng của từng thành viên; câu hỏi tương tự đặt ra cho \MAS.

\subsection{Bảo mật cấu trúc}
Topology của \MAS{} (ai nói chuyện với ai) quyết định hiệu năng và thường là kết quả của quá trình tối ưu tốn kém, nên có thể coi là intellectual property. Đồng thời, biết được cấu trúc giúp kẻ tấn công xác định key agent để tấn công có chủ đích (jailbreak, prompt injection). Các nghiên cứu gần đây cho thấy cấu trúc này có thể bị suy ra chỉ qua truy vấn hộp đen~\cite{p1mirage}.

\section{Một số kỹ thuật liên quan}

\begin{itemize}
  \item \textbf{Mixture-of-Experts (MoE):} một bộ cổng (gating) chọn một tập con expert cho từng đầu vào. Đây là nền tảng ý tưởng của DAIN (\#7), nơi các ``agent'' thực chất là các mô-đun expert có trao đổi thông điệp.
  \item \textbf{Retrieval-Augmented Generation (RAG):} bổ sung tri thức retrieved từ kho tài liệu hoặc knowledge graph vào ngữ cảnh sinh. MAC-RRG (\#2) dùng hai ``agent'' retrieval tri thức có cấu trúc và phi cấu trúc.
  \item \textbf{Học tăng cường cho tool-using agent:} sau giai đoạn SFT, mô hình được tối ưu bằng các thuật toán như GRPO với phần thưởng về kết quả, định dạng và quá trình gọi công cụ. UrbanAgent (\#6) áp dụng hướng này.
  \item \textbf{Attention steering:} điều chỉnh logit hoặc trọng số chú ý của \LLM{} tại thời điểm suy luận để nhấn mạnh một số đoạn văn bản mà không cần huấn luyện lại, ví dụ SPA hay PASTA. AGENT-RADAR (\#8) dựa trên kỹ thuật này.
\end{itemize}

\section{Các độ đo đánh giá thường gặp}
\label{sec:metrics}

Tám nghiên cứu dùng nhiều độ đo khác nhau tùy theo bài toán. Bảng~\ref{tab:metrics} tóm tắt các độ đo xuất hiện trong báo cáo để người đọc có thể diễn giải đúng các con số ở các chương sau.

\begin{table}[!htbp]
\centering
\small
\caption{Các độ đo đánh giá được sử dụng trong tám nghiên cứu.}
\label{tab:metrics}
\begin{tabularx}{\textwidth}{L{3.2cm} Y L{2.8cm}}
\toprule
\textbf{Độ đo} & \textbf{Ý nghĩa} & \textbf{Dùng trong} \\
\midrule
Accuracy, F1 & Tỉ lệ trả lời đúng; F1 cân bằng precision và recall (dùng cho QA nhiều bước) & AGENT-RADAR, ProgRouter, DAIN, HalluProp \\
AUC / AUROC & Khả năng phân biệt hai lớp, 0.5 là ngẫu nhiên, 1.0 là hoàn hảo & MIRAGE (thấp là tốt), HalluProp (cao là tốt) \\
Hit@k, Precision@k & Agent thật sự lỗi có nằm trong top-$k$ agent bị dự báo rủi ro nhất không & HalluProp \\
BLEU, ROUGE-L, METEOR, CIDEr & Mức trùng khớp n-gram giữa văn bản sinh ra và văn bản tham chiếu & MAC-RRG \\
Clinical efficacy (CE) & Precision/Recall/F1 của các nhãn bệnh lý trích từ báo cáo sinh ra & MAC-RRG \\
$R^2$, Spearman $\rho$ & Mức giải thích phương sai và tương quan hạng của dự đoán hồi quy & UrbanAgent \\
Năng lượng (J), thời gian (s) & Chi phí vận hành thực tế của hệ thống & ProgRouter, AGENT-RADAR, HalluProp \\
Tỉ lệ đa số đúng & Tỉ lệ lần chạy mà đa số agent trả lời đúng & TIDE \\
Propagation order & Số bước chuyển tiếp của một bằng chứng từ nguồn gốc đến agent tiếp nhận & TIDE \\
\bottomrule
\end{tabularx}
\end{table}

Lưu ý là cùng một độ đo có thể mang ý nghĩa ngược nhau: với MIRAGE, AUC là độ chính xác của \emph{kẻ tấn công} nên càng thấp càng tốt; với HalluProp, AUROC là độ chính xác của \emph{bộ dự báo rủi ro} nên càng cao càng tốt.

\chapter{TỔNG QUAN VÀ PHÂN LOẠI CÁC NGHIÊN CỨU}
\label{chap:taxonomy}

\section{Khung phân loại bốn nhóm}

Dựa trên câu hỏi nghiên cứu chính mà mỗi bài giải quyết, nhóm phân loại tám nghiên cứu thành bốn hướng (Hình~\ref{fig:taxonomy}):

\begin{enumerate}
  \item \textbf{Safety \& Robustness} (\textcolor{gSafe}{\textbf{S}}): hệ thống có chống chịu được tấn công hay thông tin sai lệch không? Gồm MIRAGE (\#1) -- bảo vệ bí mật topology, và TIDE (\#4) -- đánh giá ảnh hưởng của thông tin sai đến quá trình khôi phục sự thật tập thể.
  \item \textbf{Orchestration \& Control} (\textcolor{gOrch}{\textbf{O}}): quyết định ai chạy, chạy bằng mô hình nào, và agent nào có rủi ro. Gồm ProgRouter (\#3) -- model routing theo tiến độ dưới ràng buộc chi phí, và HalluProp (\#5) -- suy luận rủi ro thất bại trước khi hệ thống chạy.
  \item \textbf{Communication \& Interaction} (\textcolor{gComm}{\textbf{C}}): các agent nên trao đổi thông tin gì, với ai và ở mức độ nào. Gồm DAIN (\#7) -- mạng tương tác động giữa các agent chuyên biệt cho dữ liệu đa phương thức, và AGENT-RADAR (\#8) -- attention steering theo mức độ liên quan của ngữ cảnh.
  \item \textbf{Domain Applications} (\textcolor{gApp}{\textbf{A}}): áp dụng multi-agent collaboration vào bài toán chuyên ngành. Gồm MAC-RRG (\#2) -- radiology report generation ngực, và UrbanAgent (\#6) -- dự đoán chỉ số kinh tế -- xã hội vùng đô thị.
\end{enumerate}

\begin{figure}[!htbp]
\centering
\resizebox{\linewidth}{!}{%
\begin{tikzpicture}[
  font=\small,
  root/.style={draw,rounded corners,fill=black!8,minimum width=5.2cm,minimum height=0.9cm,align=center,font=\bfseries\small},
  grp/.style={draw,rounded corners,minimum width=3.4cm,minimum height=0.9cm,align=center,font=\bfseries\footnotesize,text=white},
  leaf/.style={draw,rounded corners=1pt,minimum width=3.4cm,align=center,font=\footnotesize,fill=white},
  arr/.style={-{Stealth[length=2mm]},thick,black!60}
]
\node[root] (r) at (0,0) {Multi-Agent Systems \&\\Collaborative Reasoning};
\node[grp,fill=gSafe] (s) at (-6.0,-1.9) {Safety \&\\Robustness};
\node[grp,fill=gOrch] (o) at (-2.0,-1.9) {Orchestration\\\& Control};
\node[grp,fill=gComm] (c) at (2.0,-1.9)  {Communication\\\& Interaction};
\node[grp,fill=gApp]  (a) at (6.0,-1.9)  {Domain\\Applications};
\foreach \g in {s,o,c,a} \draw[arr] (r.south) -- (\g.north);

\node[leaf,below=0.45cm of s] (s1) {\#1 MIRAGE\\\scriptsize ẩn topology};
\node[leaf,below=0.2cm of s1] (s2) {\#4 TIDE\\\scriptsize thông tin sai \& sự thật phân tán};
\node[leaf,below=0.45cm of o] (o1) {\#3 ProgRouter\\\scriptsize routing theo tiến độ};
\node[leaf,below=0.2cm of o1] (o2) {\#5 HalluProp\\\scriptsize rủi ro trước khi chạy};
\node[leaf,below=0.45cm of c] (c1) {\#7 DAIN\\\scriptsize mạng tương tác động};
\node[leaf,below=0.2cm of c1] (c2) {\#8 AGENT-RADAR\\\scriptsize attention steering};
\node[leaf,below=0.45cm of a] (a1) {\#2 MAC-RRG\\\scriptsize báo cáo X-quang};
\node[leaf,below=0.2cm of a1] (a2) {\#6 UrbanAgent\\\scriptsize urban region profiling};
\foreach \g/\l in {s/s1,o/o1,c/c1,a/a1} \draw[arr] (\g) -- (\l);
\end{tikzpicture}}
\caption{Khung phân loại bốn nhóm cho tám nghiên cứu được khảo sát.}
\label{fig:taxonomy}
\end{figure}

\section{Ma trận multi-label}

Phân loại một nhãn chỉ phản ánh mục tiêu \emph{chính}. Trên thực tế, hầu hết các bài đều chạm tới nhiều nhóm: HalluProp được xếp vào điều phối vì đầu ra của nó dùng để cắt tỉa cạnh hoặc sửa vai trò trước khi chạy, nhưng bản chất là phương pháp an toàn (bài được nộp vào chuyên mục cs.CR); DAIN vừa là phương pháp giao tiếp vừa là phương pháp tiết kiệm tính toán qua sparse activation. Bảng~\ref{tab:crosslabel} thể hiện quan hệ đa nhãn này.

\begin{table}[!htbp]
\centering
\small
\caption{Ma trận multi-label: \textbf{\large$\bullet$} = hướng chính, $\circ$ = hướng phụ.}
\label{tab:crosslabel}
\begin{tabular}{l C{1.8cm} C{1.8cm} C{1.8cm} C{2.25cm}}
\toprule
\textbf{Nghiên cứu} & \textcolor{gSafe}{\textbf{Safety}} & \textcolor{gOrch}{\textbf{Orchestr.}} & \textcolor{gComm}{\textbf{Comm.}} & \textcolor{gApp}{\textbf{Application}} \\
\midrule
\#1 MIRAGE      & \textbf{\large$\bullet$} &          & $\circ$  &          \\
\#4 TIDE        & \textbf{\large$\bullet$} &          & $\circ$  &          \\
\#3 ProgRouter  &          & \textbf{\large$\bullet$} &          &          \\
\#5 HalluProp   & $\circ$  & \textbf{\large$\bullet$} & $\circ$  &          \\
\#7 DAIN        &          & $\circ$  & \textbf{\large$\bullet$} & $\circ$ (y tế) \\
\#8 AGENT-RADAR &          & $\circ$  & \textbf{\large$\bullet$} &          \\
\#2 MAC-RRG     &          & $\circ$  &          & \textbf{\large$\bullet$} \\
\#6 UrbanAgent  &          &          & $\circ$  & \textbf{\large$\bullet$} \\
\bottomrule
\end{tabular}
\end{table}

\section{Vị trí can thiệp trong lifecycle của MAS}

Một cách nhìn bổ sung là xét \emph{thời điểm} mỗi phương pháp tác động lên vòng đời của một \MAS{} (Hình~\ref{fig:lifecycle}):
\begin{itemize}
  \item \textbf{Trước khi chạy (pre-hoc):} HalluProp chỉ dùng role prompt, câu hỏi và đồ thị để ước lượng rủi ro, chưa có agent nào sinh ra câu trả lời.
  \item \textbf{Trong khi chạy -- điều phối:} ProgRouter chọn mô hình cho từng bước; DAIN chọn agent nào được kích hoạt cho từng mẫu.
  \item \textbf{Trong khi chạy -- giao tiếp:} AGENT-RADAR quyết định agent chú ý vào câu nào trong lịch sử; DAIN học cường độ kết nối giữa các agent; UrbanAgent xây đồ thị quan hệ (corroboration, rectification, supervisory) giữa các agent.
  \item \textbf{Khi phơi bày kết quả:} MIRAGE tách ``execution view'' khỏi ``exposure view'' để đánh lừa kẻ tấn công.
  \item \textbf{Đánh giá sau cùng:} TIDE là một khung đánh giá, theo dõi dòng chảy bằng chứng đúng và sai qua toàn bộ quá trình.
\end{itemize}

\begin{figure}[!htbp]
\centering
\resizebox{\linewidth}{!}{%
\begin{tikzpicture}[
  font=\footnotesize,
  stage/.style={draw,thick,rounded corners,minimum width=2.75cm,minimum height=1.0cm,align=center,fill=black!5,font=\bfseries\footnotesize},
  tag/.style={draw,rounded corners=1pt,align=center,font=\scriptsize,minimum width=2.55cm,inner sep=2pt},
  arr/.style={-{Stealth[length=2.5mm]},very thick,black!55}
]
\node[stage] (p0) at (0,0)    {Thiết kế\\vai trò \& đồ thị};
\node[stage] (p1) at (3.15,0) {Điều phối\\\& kích hoạt};
\node[stage] (p2) at (6.3,0)  {Giao tiếp\\\& lập luận};
\node[stage] (p3) at (9.45,0) {Tổng hợp\\\& phơi bày};
\node[stage] (p4) at (12.6,0) {Đánh giá\\\& kiểm toán};
\draw[arr] (p0)--(p1); \draw[arr] (p1)--(p2); \draw[arr] (p2)--(p3); \draw[arr] (p3)--(p4);

\node[tag,fill=gOrch!15,below=0.4cm of p0] (h) {\#5 HalluProp\\(pre-hoc)};
\node[tag,fill=gOrch!15,below=0.4cm of p1] (pr) {\#3 ProgRouter};
\node[tag,fill=gComm!15,below=0.15cm of pr] (d1) {\#7 DAIN (cổng)};
\node[tag,fill=gComm!15,below=0.4cm of p2] (ra) {\#8 AGENT-RADAR};
\node[tag,fill=gComm!15,below=0.15cm of ra] (d2) {\#7 DAIN (đồ thị)};
\node[tag,fill=gApp!20,below=0.15cm of d2] (u) {\#6 UrbanAgent};
\node[tag,fill=gApp!20,below=0.15cm of u] (m) {\#2 MAC-RRG};
\node[tag,fill=gSafe!15,below=0.4cm of p3] (mi) {\#1 MIRAGE};
\node[tag,fill=gSafe!15,below=0.4cm of p4] (t) {\#4 TIDE};
\end{tikzpicture}}
\caption{Vị trí can thiệp của từng nghiên cứu trong vòng đời một multi-agent system.}
\label{fig:lifecycle}
\end{figure}

\section{Bản chất của ``agent'' trong các nghiên cứu}

Một quan sát quan trọng khi đọc kỹ tám bài là thuật ngữ ``agent'' được dùng với nghĩa rất khác nhau (Bảng~\ref{tab:agentnature}). Chỉ năm bài (\#1, \#3, \#4, \#5, \#8) dùng \emph{LLM agent thực sự} trao đổi bằng ngôn ngữ tự nhiên. UrbanAgent dùng LLM agent ở giai đoạn đầu nhưng giai đoạn ``debate'' là một bộ tổng hợp nơ-ron được huấn luyện. DAIN gọi các mô-đun nơ-ron con là ``agent'', về bản chất là một MoE có trao đổi thông điệp giữa các expert. MAC-RRG gọi các mô-đun retrieval tất định là ``agent''. Sự khác biệt này ảnh hưởng trực tiếp đến khả năng so sánh và được thảo luận thêm ở Chương~\ref{chap:comparison}.

\begin{table}[!htbp]
\centering
\small
\caption{Bản chất của agent và phương thức giao tiếp trong từng nghiên cứu.}
\label{tab:agentnature}
\begin{tabularx}{\textwidth}{L{2.4cm} L{3.3cm} Y L{2.4cm}}
\toprule
\textbf{Nghiên cứu} & \textbf{Agent là gì} & \textbf{Phương thức giao tiếp} & \textbf{Cần huấn luyện?} \\
\midrule
\#1 MIRAGE & LLM theo vai trò & Văn bản trên đồ thị DAG & Không \\
\#3 ProgRouter & LLM (vai trò thực thi) & Coordinator giao việc qua sổ cái (ledger) & Bộ dự báo học trực tuyến \\
\#4 TIDE & LLM (5 agent đồng nhất) & Phát quảng bá công khai + kênh riêng, 3 vòng & Không \\
\#5 HalluProp & LLM theo vai trò & Văn bản trên đồ thị, 3 vòng & Risk predictor có giám sát \\
\#8 AGENT-RADAR & LLM theo vai trò & Văn bản, đồ thị ngẫu nhiên/phân lớp/đầy đủ/debate & Không \\
\#6 UrbanAgent & VLM dùng công cụ (mỗi phương thức một agent) & Đồ thị quan hệ học được (truyền vector) & Có (SFT + GRPO + GNN) \\
\#7 DAIN & Mô-đun nơ-ron & Vector qua low-rank bottleneck & Có (đầu-cuối) \\
\#2 MAC-RRG & Mô-đun retrieval tất định & Ghép token tri thức vào tiền tố LLM & Có (LLM sinh báo cáo) \\
\bottomrule
\end{tabularx}
\end{table}

\chapter{TRÌNH BÀY CÁC NGHIÊN CỨU TIÊU BIỂU}
\label{chap:papers}

Chương này trình bày chi tiết tám nghiên cứu theo bốn nhóm của khung phân loại. Mỗi nghiên cứu gồm: bài toán và động lực, phương pháp, thiết lập thực nghiệm, kết quả chính và nhận xét phản biện của nhóm.

% =====================================================================
\section{Nhóm Safety \& Robustness}
% =====================================================================

\subsection{MIRAGE: Che giấu topology của MAS bằng phantom topology}
\label{sec:mirage}

\noindent\textit{He, Wen, Su et al. (NTU, BUPT, PolyU), arXiv 2609.37567, 09/2026, preprint~\cite{p1mirage}.}

\paragraph{Bài toán và động lực.} Topology của \MAS{} -- đồ thị có hướng mô tả agent nào chuyển đầu ra cho agent nào -- quyết định hiệu năng lập luận và thường được tạo bởi các bộ tối ưu tốn kém như G-Designer, AGP, ARG-Designer. Tấn công CIA (Communication Inference Attack) đã chỉ ra rằng kẻ tấn công chỉ với quyền truy vấn hộp đen có thể khiến \MAS{} để lộ đầu ra trung gian của các agent, rồi suy ra cạnh dựa trên quan sát rằng hai agent được nối trực tiếp có đầu ra phụ thuộc ngữ nghĩa mạnh hơn. Việc lộ cấu trúc gây hai rủi ro: (a) lộ điểm yếu -- biết key agent để tấn công có chủ đích; (b) đánh cắp intellectual property. Các cơ chế phòng thủ \MAS{} hiện có (BlindGuard, GUARDIAN, PsySafe,~\ldots) bảo vệ \emph{nội dung} hoặc \emph{hành vi}, chưa có cơ chế nào bảo vệ \emph{tính bí mật của cấu trúc}. Mục tiêu hình thức là tìm hệ được bảo vệ $\tilde S$ sao cho
\begin{equation}
  \max_{\tilde S}\ \Delta_{topo}(\hat G, G) \quad \text{s.t.}\quad U(\tilde S) \ge U(S) - \epsilon,
\end{equation}
trong đó $\hat G$ là cấu trúc kẻ tấn công suy ra, $U$ là độ chính xác tác vụ và $\epsilon$ là mức suy giảm cho phép.

\paragraph{Phương pháp.} Ý tưởng cốt lõi là \emph{chạy} tác vụ trên đồ thị thật $G$ nhưng \emph{cho kẻ tấn công thấy} bằng chứng chỉ về một phantom graph $G'$. MIRAGE gồm ba giai đoạn:
\begin{enumerate}
  \item \textbf{Tổng hợp phantom topology.} Giữ nguyên tập agent, chỉ đi dây lại các cạnh trong không gian cạnh tôn trọng thứ tự thực thi (để không tạo chu trình). Độ lệch cấu trúc được đo bằng Structural Camouflage Ratio $\rho(H,G) = 1 - \frac{|E_H\cap E|}{|E_H\cup E|}$ (một trừ hệ số Jaccard). Lấy cảm hứng từ $k$-anonymity, topology anonymity set $K_{\rho_0}(G) = \{H : \rho(H,G)\ge\rho_0\}$ chứa các đồ thị đủ khác biệt; một phần tử được chọn làm $G'$.
  \item \textbf{Semantic Edge Realization.} Với mỗi \emph{phantom edge}, một \LLM{} sinh ``bằng chứng phụ thuộc'' giả, tuân thủ nhất quán vai trò, liên quan đến tác vụ và hợp lý về phụ thuộc. Với mỗi \emph{genuine edge cần giấu}, đầu ra được diễn đạt lại thành $M$ ứng viên và chọn ứng viên cân bằng giữa sai lệch ngữ nghĩa và độ phụ thuộc vào các agent tiền nhiệm bị giấu:
  \begin{equation}
    \tilde r_i = \arg\min_{\hat r \in R_i}\Big[D_{sem}(\hat r, r_i) + \lambda \sum_{a_j \in P_i^h} S_{dep}(r_j, \hat r)\Big].
  \end{equation}
  \item \textbf{Dual-view execution.} Execution view vẫn chạy trên $G$ nên tính toán không đổi; exposure view gồm vết lập luận đã che giấu cộng bằng chứng ảo, dùng để tạo phản hồi cho người dùng bên ngoài.
\end{enumerate}

\paragraph{Thực nghiệm.} Bốn tập dữ liệu (MMLU, GSM8K, SVAMP, HumanEval; 100 tác vụ mỗi tập), ba bộ tạo cấu trúc (G-Designer, AGP, ARG-Designer), tấn công CIA. Đối chứng: không phòng thủ, chỉ dẫn phòng ngừa (Instruction), dấu phân cách (Delimiters) và phát hiện bằng perplexity. Độ đo: AUC/ACC/F1 của cạnh suy ra (càng thấp càng tốt, AUC $\approx 0.5$ là lý tưởng).

\begin{table}[!htbp]
\centering
\small
\caption{Kết quả MIRAGE với G-Designer (AUC / ACC / F1 của tấn công, thấp hơn là tốt hơn)~\cite{p1mirage}.}
\label{tab:mirage}
\begin{tabular}{lccc}
\toprule
\textbf{Tập dữ liệu} & \textbf{Instruction} & \textbf{Delimiters} & \textbf{MIRAGE} \\
\midrule
MMLU      & 0.79 / 0.59 / 0.62 & 0.76 / 0.61 / 0.64 & \textbf{0.67 / 0.50 / 0.49} \\
GSM8K     & 0.81 / 0.73 / 0.64 & 0.80 / 0.76 / 0.67 & \textbf{0.64 / 0.51 / 0.42} \\
SVAMP     & 0.72 / 0.74 / 0.69 & 0.70 / 0.68 / 0.71 & \textbf{0.64 / 0.48 / 0.39} \\
HumanEval & 0.78 / 0.62 / 0.59 & 0.76 / 0.65 / 0.56 & \textbf{0.68 / 0.44} / 0.52 \\
\bottomrule
\end{tabular}
\end{table}

Phát hiện bằng perplexity chỉ đạt AUC 0.51--0.60 trong việc phân biệt truy vấn tấn công với truy vấn bình thường, tức gần như ngẫu nhiên. Theo tác giả, độ chính xác tác vụ ``được bảo toàn phần lớn''. Thí nghiệm loại bỏ cho thấy bỏ thành phần sinh phantom edge hoặc che giấu cạnh thật đều làm tấn công mạnh lên.

\begin{nhanxet}
\textbf{Điểm mạnh:} mở ra một chiều an toàn mới (bí mật cấu trúc); thiết kế tách thực thi/phơi bày gọn gàng; hình thức hóa theo $k$-anonymity với tham số điều chỉnh $\rho_0$. \textbf{Hạn chế:} chỉ đánh giá với một tấn công (CIA), không có kẻ tấn công thích nghi biết về MIRAGE (ví dụ phát hiện văn bản do LLM sinh); AUC sau phòng thủ vẫn khoảng 0.64--0.68, xa mức 0.5 nên việc che giấu mới chỉ một phần; quy mô nhỏ (100 tác vụ/tập); không báo cáo chi phí phát sinh do các lần gọi LLM thêm để diễn đạt lại và sinh bằng chứng; không thảo luận phương án đơn giản là không phơi bày vết lập luận trung gian.
\end{nhanxet}

% ---------------------------------------------------------------------
\subsection{TIDE: Thông tin sai làm chệch hướng khôi phục sự thật tập thể}
\label{sec:tide}

\noindent\textit{Yan, Yue, Zhao et al. (CASIA và cộng sự), arXiv 2608.03421, 08/2026, preprint~\cite{p4tide}.}

\paragraph{Bài toán và động lực.} \MAS{} thường được đánh giá qua độ chính xác câu trả lời cuối, nhưng điều đó không cho biết bằng chứng \emph{riêng} của từng agent có thực sự được đưa vào và tích hợp đúng trong thảo luận hay không. Lấy cảm hứng từ nghiên cứu ``hidden profile'' ở nhóm người~\cite{stasser1985}, bài báo đặt câu hỏi: khi sự thật bị phân tán giữa các agent và có một deceiver, bằng chứng đúng và sai cạnh tranh, lan truyền ra sao?

\paragraph{Phương pháp: khung đánh giá ForesightSafety-TIDE.} Đây là \emph{khung đánh giá có kiểm soát}, không phải phương pháp \MAS{} mới.
\begin{itemize}
  \item \textbf{Môi trường:} 120 môi trường được xây từ truyện Hi-ToM. Mỗi môi trường có một đồ vật di chuyển qua chuỗi vị trí $l_0 \to l_1 \to \cdots \to l_m$, và 5 agent, mỗi agent chỉ quan sát một phần chuỗi. Thiết kế đảm bảo hình thức hai tính chất: \emph{tập thể đủ thông tin} -- gộp bằng chứng của 4 agent không then chốt thì xác định được duy nhất $l_m$; \emph{cá nhân thiếu thông tin} -- không agent nào tự xác định được.
  \item \textbf{Điều kiện ghép cặp:} C0 -- tất cả trung thực; C1 -- key agent (người thấy bước di chuyển cuối) nói dối hướng tới decoy $d = l_{m-1}$ bằng ba cơ chế: \emph{Event} (bịa ra bước di chuyển ngược), \emph{Coverage} (kéo dài cửa sổ quan sát), \emph{Source} (gán sai thời điểm của một báo cáo cũ).
  \item \textbf{Giao thức:} phán đoán riêng (T0), cập nhật sau testimony card bắt buộc có kiểu (T1), ba vòng thảo luận công khai (R1--R3), bỏ phiếu riêng sau mỗi vòng. Mô hình: GPT-5.5, DeepSeek-v4-pro, Grok-4.5, mỗi mô hình là một hệ đồng nhất 5 agent.
  \item \textbf{Độ đo quá trình:} tỉ lệ trích dẫn và chấp nhận testimony, \emph{evidence lineage} (gốc đúng / gốc sai) và \emph{propagation order} (số bước chuyển tiếp).
\end{itemize}

\paragraph{Kết quả chính.}
\begin{itemize}
  \item Tỉ lệ đa số đúng giảm từ \textbf{72.50\%} (C0) xuống \textbf{14.17\%} (C1), giảm 58.33 điểm phần trăm; phiếu đúng cá nhân giảm từ 75.00\% xuống 24.38\%; $p < .001$ với cả ba mô hình.
  \item Trong 309 lần thất bại ở C1, khoảng một nửa (156) là đa số theo decoy và một nửa (152) là không đồng thuận.
  \item \textbf{Adoption asymmetry:} testimony sai được chấp nhận 70.7\% trong khi testimony đúng chỉ 25.4\%. Bằng chứng sai lan xa hơn: propagation order trung bình ở R3 là 2.66 so với 1.33 của bằng chứng đúng.
  \item Khoảng 85\% khoảng cách cuối cùng đã hình thành ngay sau khi chia sẻ testimony (T1). Thảo luận \emph{khuếch đại} chiều hướng bằng chứng sẵn có chứ không sửa sai.
  \item Khi deceiver rời nhóm sau khi khai, tỉ lệ đa số đúng vẫn thấp hơn 46.67 điểm so với khi agent đó rời đi im lặng: thông tin sai tồn tại độc lập với người phát tán.
  \item Thêm ``observer'' không có bằng chứng trực tiếp giảm khoảng 90\% số đa số theo decoy, nhưng \emph{không} tăng đa số đúng: họ làm nhóm thận trọng hơn chứ không sửa được lỗi.
\end{itemize}

\begin{nhanxet}
\textbf{Điểm mạnh:} thiết kế ghép cặp cô lập chính xác tác động của lừa dối; tính đủ/thiếu thông tin được kiểm chứng hình thức; truy vết evidence lineage vượt xa các độ đo chỉ nhìn kết quả; kiểm định thống kê chặt (bootstrap, McNemar). \textbf{Hạn chế:} chỉ một tác vụ hẹp (theo dõi vị trí đồ vật); mỗi ô chạy một lần; chỉ hệ đồng nhất 5 agent, kết nối đầy đủ; thí nghiệm bổ sung chỉ dùng GPT-5.5; decoy luôn là $l_{m-1}$ -- vị trí mà một số agent trung thực cũng thấy là ``lần cuối'', có thể thổi phồng sức mạnh của testimony sai; agent được báo trước ``tối đa một người không đáng tin'', một giả định mạnh; chưa đề xuất phương pháp phòng thủ.
\end{nhanxet}


% ---------------------------------------------------------------------
\subsection{So sánh trong nhóm: MIRAGE và TIDE}
\label{sec:cmp-safety}

Hai nghiên cứu cùng thuộc hướng an toàn nhưng nhìn rủi ro từ hai phía đối lập: MIRAGE xem mối đe doạ là một \emph{kẻ tấn công bên ngoài} muốn suy ra cấu trúc hệ thống, còn TIDE xem mối đe doạ là \emph{một agent bên trong} cố tình đưa thông tin sai. Bảng~\ref{tab:cmp-g1} so sánh chi tiết.

\begin{table}[!htbp]
\centering
\small
\caption{So sánh MIRAGE và TIDE.}
\label{tab:cmp-g1}
\begin{tabularx}{\textwidth}{L{3.0cm} Y Y}
\toprule
\textbf{Tiêu chí} & \textbf{MIRAGE} & \textbf{TIDE} \\
\midrule
Loại đóng góp & Phương pháp phòng thủ & Khung đánh giá (benchmark) \\
Nguồn đe doạ & Kẻ tấn công bên ngoài, truy vấn hộp đen & Agent bên trong nói dối \\
Tài sản cần bảo vệ & Tính bí mật của topology & Tính đúng đắn của kết luận tập thể \\
Topology được xét & DAG tối ưu (G-Designer, AGP, ARG-Designer) & Kết nối đầy đủ, phát quảng bá + kênh riêng \\
Backbone & Không nêu trong thân bài & GPT-5.5, DeepSeek-v4-pro, Grok-4.5 \\
Độ đo chính & AUC/ACC/F1 của cạnh bị suy ra & Đa số đúng/decoy, evidence lineage, propagation order \\
Kết quả nổi bật & AUC tấn công giảm còn 0.64--0.68 & Đa số đúng giảm 72.50\% $\to$ 14.17\% \\
Kiểm định thống kê & Không & Có (bootstrap, McNemar) \\
Hạn chế chính & Một loại tấn công (CIA); chưa có kẻ tấn công thích nghi & Một tác vụ (Hi-ToM); chưa đề xuất phòng thủ \\
\bottomrule
\end{tabularx}
\end{table}

Hai bài bổ sung cho nhau theo một logic rõ ràng. TIDE cho thấy một khi thông tin sai đã vào hệ thống thì gần như không thể sửa bằng thảo luận; MIRAGE cho thấy có thể làm kẻ tấn công khó xác định agent then chốt -- tức khó chọn đúng chỗ để đưa thông tin sai vào. Tuy vậy, chưa bài nào đề xuất cơ chế \emph{khôi phục} sự thật khi tấn công đã xảy ra; đây là khoảng trống được bàn ở Chương~\ref{chap:gaps}.

% =====================================================================
\section{Nhóm Orchestration \& Control}
% =====================================================================

\subsection{ProgRouter: Online progress-guided orchestration}
\label{sec:progrouter}

\noindent\textit{Li, Abdelmoniem, Wang (Aston, QMUL, Exeter), arXiv 2608.25992, 08/2026, preprint~\cite{p3progrouter}.}

\paragraph{Bài toán và động lực.} Quy trình multi-agent (kiểu Magentic-One~\cite{magenticone}) gọi \LLM{} nhiều lần với ngữ cảnh dài, rất tốn kém. Các router hiện có (RouteLLM, CASCADIA, MasRouter) ra quyết định \emph{một lần cho mỗi câu hỏi}, trong khi quy trình multi-agent cần \emph{chuỗi} quyết định phụ thuộc trạng thái: ở mỗi bước, vai trò được gọi nên dùng mô hình nào? Ba khó khăn: quy trình mở và có trạng thái; tiến độ phải ước lượng từ trạng thái trung gian nhiễu; ngân sách phải được tuân thủ khi chưa biết các bước tương lai.

\paragraph{Phương pháp.} Một coordinator agent $C$ lập kế hoạch và chọn \emph{vai trò} tiếp theo; ProgRouter chọn \emph{mô hình} cho vai trò đó từ model zoo. Bài toán là cực đại hóa chất lượng trung bình dưới ràng buộc chi phí (năng lượng GPU, Jun) và thời gian theo từng tác vụ cũng như trung bình dài hạn. Ba thành phần:
\begin{enumerate}
  \item \textbf{Multi-view progress scorer} $g(s_t)\in[0,1]$, kết hợp phân loại trạng thái tổng thể $c_t$ (không hợp lệ / có thể cứu / thành công một phần / hoàn tất), tỉ lệ yêu cầu con đã hoàn thành $r_t$, xu hướng tiến độ $d_t$ và chất lượng trạng thái $e_t$:
  \begin{equation}
    g(s_t) = b(c_t) + \alpha_{c_t} r_t + \beta_{c_t} d_t + \gamma_{c_t} e_t.
  \end{equation}
  \item \textbf{Dual-path progress predictor}: nhánh cấu trúc (cây hồi quy trên đặc trưng bảng) và nhánh ngữ nghĩa (cây hồi quy trên embedding MiniLM của bản tóm tắt trạng thái), kết hợp bằng một meta-learner XGBoost, dự báo mức tăng tiến độ $\hat y_t$ nếu dùng mô hình $m_t$.
  \item \textbf{Online decision-making theo Lyapunov drift-plus-penalty} với virtual cost queue $Q_w$:
  \begin{equation}
    \text{score}(m_t) = V(1-g(s_t))\,\hat y_t(m_t) - Q_w\big(E(m_t)-\tilde E\big) - c^{\Gamma}_t\,\Gamma(m_t) - c^{E}_t\,E(m_t).
  \end{equation}
  Mức tăng tiến độ có giá trị hơn khi tác vụ còn xa đích; các hệ số phạt tăng theo hàm mũ khi ngân sách sắp cạn. Bộ dự báo được học trực tuyến với khám phá $\varepsilon$-greedy, không cần dữ liệu huấn luyện ngoại tuyến.
\end{enumerate}

\paragraph{Thực nghiệm.} Khung AG2/AutoGen; coordinator Qwen3-Coder-Next 80B; model zoo gồm các mô hình Qwen, Granite, Gemma từ 0.5B đến 35B. Bốn bộ dữ liệu: HumanEval+, MBPP, MATH-500, ASQA. Đối chứng: MasRouter, CASCADIA, Educated Guessing, các chính sách một mô hình cố định. Năng lượng đo thực bằng NVML.

\begin{table}[!htbp]
\centering
\small
\caption{Kết quả ProgRouter (độ chính xác \% / năng lượng J). $^{\dagger}$ = vượt ngân sách năng lượng~\cite{p3progrouter}.}
\label{tab:progrouter}
\begin{tabular}{lcccc}
\toprule
\textbf{Phương pháp} & \textbf{HumanEval+} & \textbf{MBPP} & \textbf{MATH-500} & \textbf{ASQA} \\
\midrule
Ngân sách $\tilde E$ & 4\,800 & 4\,500 & 7\,000 & 19\,000 \\
\midrule
MasRouter         & 90.9 / 4\,483 & 67.8 / 3\,700 & 73.6 / 8\,294$^{\dagger}$ & 89.8 / 16\,368 \\
CASCADIA          & 84.8 / 4\,658 & 78.5 / 3\,857 & \textbf{87.8} / 6\,875 & 89.8 / 27\,857$^{\dagger}$ \\
Educated Guessing & 91.5 / 4\,916$^{\dagger}$ & 65.2 / 3\,831 & 79.3 / 7\,023$^{\dagger}$ & 90.8 / 23\,423$^{\dagger}$ \\
\textbf{ProgRouter} & \textbf{93.0} / 4\,796 & \textbf{79.4 / 3\,376} & 84.3 / \textbf{6\,112} & \textbf{92.1} / 18\,373 \\
\bottomrule
\end{tabular}
\end{table}

Thí nghiệm loại bỏ trên HumanEval+: bỏ bộ dự báo còn 89.0\%, bỏ hệ số phạt ngân sách còn 87.8\%, và một chiến lược tham lam ``tiến độ trên chi phí'' ngây thơ chỉ đạt 17.7\% với 7\,797 J. Điều này cho thấy chính cơ chế ngân sách mới làm nên hiệu quả. Meta-learner có sai số MAE thấp nhất (0.0720) so với từng nhánh riêng lẻ.

\begin{nhanxet}
\textbf{Điểm mạnh:} routing theo từng bước, có nhận thức trạng thái; không cần dữ liệu ngoại tuyến; cơ chế ngân sách có nền tảng lý thuyết; đo năng lượng GPU thực thay cho token. \textbf{Hạn chế:} trên MATH-500 thua CASCADIA 3.5 điểm độ chính xác; tuyên bố ``cao nhất'' trên ASQA chỉ đúng trong số các phương pháp không vượt ngân sách (Qwen3.6-35B đơn lẻ đạt 92.3); kết quả chỉ báo cáo sau giai đoạn khởi động, bỏ qua chi phí khởi động; thí nghiệm loại bỏ chỉ trên một tập; không báo cáo phương sai; progress scorer cần định nghĩa thủ công các mốc cho từng miền.
\end{nhanxet}

% ---------------------------------------------------------------------
\subsection{HalluProp: Suy luận rủi ro thất bại trước khi agent lên tiếng}
\label{sec:halluprop}

\noindent\textit{Lin, Wang, Qian et al. (ĐH Công Thương Chiết Giang, ĐH Chiết Giang), arXiv 2607.26836, 07/2026, preprint~\cite{p5halluprop}.}

\paragraph{Bài toán và động lực.} Hallucination của một agent có thể lan theo các communication edge, được agent khác tiếp nhận và khuếch đại thành cascading failure. Các phương pháp phát hiện hiện tại đều \emph{post-hoc}: kiểm tra từng câu trả lời (SelfCheckGPT) hoặc quy trách nhiệm dựa trên đồ thị tương tác (GUARDIAN) -- đều cần chạy hết hệ thống và chậm. Câu hỏi đặt ra: có thể suy ra rủi ro của từng agent \emph{trước khi} có bất kỳ tương tác nào, chỉ từ role prompt, câu hỏi và topology không?

\paragraph{Phương pháp.} Rủi ro được tách thành \emph{intrinsic risk} của từng đỉnh $\pi_i$ và \emph{propagation risk} của từng cạnh $t_{j\to i}$:
\begin{itemize}
  \item \textbf{Intrinsic risk:} mã hóa vai trò $e^R_i$ và câu hỏi $e^Q$ bằng bộ mã hóa văn bản đóng băng (BERT). Một cổng phụ thuộc câu hỏi tách vai trò thành phần \emph{năng lực} khớp với câu hỏi $\kappa_i$ và phần \emph{dễ tổn thương} còn lại $\rho_i$; $\pi_i = \sigma(w_\rho^\top\rho_i - w_\kappa^\top\kappa_i + b)$.
  \item \textbf{Propagation risk:} tích của năng lực ảnh hưởng của bên gửi và độ dễ tổn thương của bên nhận, nhân với hệ số cấu trúc dựa trên bậc vào/ra (luôn $\ge 1$, tức cấu trúc chỉ khuếch đại rủi ro).
  \item \textbf{Suy luận cấp hệ thống bằng Noisy-OR khả vi}, lặp điểm bất động $K$ bước (tốt nhất $K=3$), xử lý được cả đồ thị có chu trình:
  \begin{equation}
    h_i^{(k+1)} = 1 - (1-\pi_i)\prod_{j\in\mathcal N_{in}(i)}\big(1 - t_{j\to i}\, h_j^{(k)}\big),\qquad h^{(0)} = \pi.
  \end{equation}
  \item \textbf{Can thiệp:} với agent rủi ro cao nhất, hoặc \emph{cắt} các cạnh ra của nó, hoặc \emph{tinh chỉnh} prompt thành vai trò ``kiểm chứng thận trọng''.
\end{itemize}

\paragraph{Thực nghiệm.} MMLU, MATH, MedQA (500 câu mỗi tập); bốn backbone GPT-3.5, GPT-4o, Llama-3-8B, Qwen-2-7B; đồ thị ngẫu nhiên với 3--7 agent, 3 vòng; vai trò được nhiễu có chủ đích (lệch miền, mơ hồ, thiên lệch). Nhãn ``hallucination'' là việc đầu ra của agent sai so với đáp án.

\paragraph{Kết quả chính.} AUROC trung bình 84.6 (cao nhất 91.5 với GPT-4o trên MMLU; MedQA thấp hơn khoảng 10 điểm). Thời gian chẩn đoán dưới 1 giây so với 53--308 giây của các phương pháp post-hoc (nhanh hơn trên 65 lần). Sau can thiệp, độ chính xác trung bình là 67.6\% (tinh chỉnh) và 63.6\% (cắt cạnh) so với 63.5\% của GUARDIAN. Loại bỏ mô hình hóa intrinsic risk, đặc trưng cấu trúc hay lan truyền làm AUROC giảm lần lượt khoảng 28\%, 15\% và 21\%. Kết quả ổn định trên các cấu trúc hình sao, chuỗi và phân cấp (chênh lệch trong khoảng 2\%).

\begin{nhanxet}
\textbf{Điểm mạnh:} hướng tiếp cận \emph{pre-hoc} mới; phân rã đỉnh/cạnh dễ giải thích; chẩn đoán dưới một giây, không cần gọi LLM; có lập luận hội tụ. \textbf{Hạn chế:} ``hallucination'' thực chất là \emph{trả lời sai} trên câu hỏi trắc nghiệm/toán; tín hiệu có thể đến chủ yếu từ các vai trò bị nhiễu có chủ đích -- tính thực tế chưa được chứng minh; là phương pháp có giám sát, cần nhãn từ các lần chạy thật, chưa kiểm tra chuyển giao sang tác vụ/mô hình mới; đồ thị nhỏ ($\le 7$ agent); so sánh tốc độ bỏ qua chi phí huấn luyện; số liệu mô hình đầy đủ chênh nhẹ giữa hai bảng (84.6 và 84.8).
\end{nhanxet}


% ---------------------------------------------------------------------
\subsection{So sánh trong nhóm: ProgRouter và HalluProp}
\label{sec:cmp-orch}

Cả hai đều ra quyết định \emph{điều khiển} hệ thống nhưng ở hai thời điểm và với hai mục tiêu khác nhau (Bảng~\ref{tab:cmp-g2}).

\begin{table}[!htbp]
\centering
\small
\caption{So sánh ProgRouter và HalluProp.}
\label{tab:cmp-g2}
\begin{tabularx}{\textwidth}{L{3.0cm} Y Y}
\toprule
\textbf{Tiêu chí} & \textbf{ProgRouter} & \textbf{HalluProp} \\
\midrule
Thời điểm can thiệp & Trong khi chạy, từng bước & Trước khi chạy (pre-hoc) \\
Câu hỏi quyết định & Vai trò này, bước này dùng model nào? & Agent/cạnh nào có rủi ro cao? \\
Mục tiêu tối ưu & Chất lượng trong ngân sách năng lượng/thời gian & Độ tin cậy (giảm lỗi lan truyền) \\
Tín hiệu đầu vào & Trạng thái workflow (ledger), lịch sử dispatch & Role prompt, câu hỏi, topology \\
Cách học & Học trực tuyến, không cần dữ liệu ngoại tuyến & Có giám sát, cần nhãn từ các lần chạy \\
Mô hình bên trong & Cây hồi quy + XGBoost & BERT đóng băng + Noisy-OR khả vi \\
Chi phí của chính phương pháp & Rất nhẹ (cây hồi quy) & $<$1 giây/chẩn đoán (chưa tính huấn luyện) \\
Tương thích model API & Có & Có \\
Kết quả nổi bật & 93.0\% HumanEval+ trong ngân sách & AUROC 84.6; accuracy 67.6\% sau refinement \\
\bottomrule
\end{tabularx}
\end{table}

Điểm chung đáng chú ý là cả hai đều chuyển từ quyết định \emph{tĩnh} sang quyết định \emph{phụ thuộc ngữ cảnh}: ProgRouter phụ thuộc tiến độ hiện tại, HalluProp phụ thuộc câu hỏi cụ thể. Hai phương pháp cũng kết hợp tự nhiên: rủi ro do HalluProp dự báo có thể trở thành một khoản phạt trong hàm điểm của ProgRouter, để các bước rủi ro cao được giao cho model mạnh hơn (đề xuất ở Chương~\ref{chap:gaps}).

% =====================================================================
\section{Nhóm Communication \& Interaction}
% =====================================================================

\subsection{DAIN: Dynamic Agent-based Interaction Network cho suy luận đa phương thức}
\label{sec:dain}

\noindent\textit{Chen, Li, Wang et al. (Viện Hàn lâm Khoa học Trung Quốc), arXiv 2606.30189, 06/2026, preprint~\cite{p7dain}.}

\paragraph{Bài toán và động lực.} Kết hợp dữ liệu đa phương thức (ảnh, văn bản, âm thanh, hồ sơ lâm sàng) là bài toán cốt lõi, nhất là trong chẩn đoán y tế. Partial Information Decomposition (PID) chia thông tin đa phương thức thành \emph{uniqueness}, \emph{redundancy} và \emph{synergy}, nhưng hiếm khi được tích hợp vào mô hình đầu-cuối. Các phương pháp Mixture-of-Experts (MoE) đa phương thức dùng cấu trúc cộng tác \emph{tĩnh}, không chọn expert theo ngữ cảnh và không cho các expert trao đổi thông tin.

\paragraph{Phương pháp.} DAIN coi việc kết hợp đa phương thức là quá trình ra quyết định cộng tác của $K=8$ ``interaction agent'' -- mỗi agent là một \emph{mô-đun nơ-ron}, không phải LLM -- gồm 3 synergy agent, 3 uniqueness agent và 2 redundancy agent.
\begin{enumerate}
  \item \textbf{Encoding \& context analysis:} mỗi phương thức qua bộ mã hóa riêng (ResNet-50, BERT-base), ghép thành vector ngữ cảnh $e_c$.
  \item \textbf{Dynamic scheduling bởi Meta-Controller:} cổng thưa chọn agent nào được kích hoạt, đồng thời sinh soft communication graph phụ thuộc đầu vào:
  \begin{equation}
    g = \mathrm{Top}_\tau\big(\mathrm{Softmax}(G_{gate}(e_c))\big),\qquad \Phi = \sigma\big(G_{comm}(e_c)\big)\in\mathbb{R}^{K\times K}.
  \end{equation}
  \item \textbf{Gated execution \& compressed communication:} thông điệp $m_{i\leftarrow j} = \Phi_{ij}\cdot W_B h_j$ qua low-rank bottleneck; một vòng cập nhật trạng thái.
  \item \textbf{Consensus fusion:} chú ý đa đầu trên trạng thái các agent được kích hoạt.
\end{enumerate}
Hàm mục tiêu kết hợp mất mát tác vụ, specialization loss dựa trên nhiễu loạn và hai chính quy hóa khuyến khích sparse activation và giao tiếp thưa:
\begin{equation}
  \mathcal L = \mathcal L_{task} + \alpha\,\mathcal L_{spec} + \lambda_{act}\|g\|_1 + \lambda_{comm}\|\Phi\|_F^2.
\end{equation}

\paragraph{Thực nghiệm và kết quả.} Năm tập dữ liệu: ADNI (Alzheimer), MIMIC-IV (tử vong ICU), MM-IMDB (thể loại phim), CMU-MOSI (cảm xúc, MSE), ENRICO (giao diện di động). DAIN vượt các đối chứng (Early/Late Fusion, MoE, Sparse MoE, MMoE) trên cả năm tập, ví dụ ADNI 73.5 so với 70.9, MOSI MSE 2.41 so với 2.48. Trung bình chỉ 3.2--5.5 trên 8 agent được kích hoạt, mật độ đồ thị 0.19--0.42. Ngân sách agent tối ưu là $\tau = 4$. Bỏ giao tiếp làm giảm 1.4 điểm trên ADNI; bỏ dynamic scheduling (kích hoạt tất cả) giảm 1.7 điểm. Phân tích kích hoạt cho thấy bệnh nhân bình thường dùng chủ yếu uniqueness agent, còn ca suy giảm nhận thức dùng nhiều synergy agent hơn.

\begin{nhanxet}
\textbf{Điểm mạnh:} thiết kế ba bước (lập lịch -- giao tiếp -- đồng thuận) rõ ràng, gắn với lý thuyết PID; sparse activation có số liệu theo từng tập; có khả năng diễn giải qua hồ sơ kích hoạt. \textbf{Hạn chế:} ``agent'' chỉ là ẩn dụ -- hệ thống thực chất là MoE có cổng và truyền thông điệp giữa expert; số liệu mức cải thiện không nhất quán giữa Bảng 1, Bảng 3 và văn bản; phép cộng số tham số ``hiệu dụng'' không khớp với con số 5.8M được nêu; tuyên bố báo cáo độ lệch chuẩn nhưng bảng không có; đối chứng yếu và cũ; nhiều trích dẫn lạc đề.
\end{nhanxet}

% ---------------------------------------------------------------------
\subsection{AGENT-RADAR: Attention steering theo context relevance}
\label{sec:radar}

\noindent\textit{Zhang, Tian, Zhang (Purdue University), arXiv 2605.30136, 05/2026, preprint~\cite{p8radar}.}

\paragraph{Bài toán và động lực.} Lịch sử hội thoại trong \MAS{} tăng nhanh và chứa nhiều nội dung lặp lại, khiến mô tả nhiệm vụ gốc và bằng chứng then chốt bị chôn vùi giữa ngữ cảnh (``lost in the middle''~\cite{lostmiddle}). Các giải pháp tóm tắt (compaction) hoặc cắt tỉa (AgentPrune, AgentDropout) đều làm mất thông tin, có thể loại bỏ đúng những tín hiệu tinh tế nhưng quan trọng. Đề xuất: \emph{giữ nguyên toàn bộ} lịch sử và dùng \emph{soft attention steering} hướng về nội dung liên quan.

\paragraph{Phương pháp.} AGENT-RADAR là mô-đun cắm thêm, không cần huấn luyện. Với agent $v_i$ ở vòng $t$ và câu hỏi hiện tại $q^{(t)}$, mỗi câu $s_k$ trong thông điệp $o_j^{(\tau)}$ của agent $v_j$ ở vòng $\tau$ được chấm điểm bằng tích ba thành phần:
\begin{equation}
  \text{score}(s_k) = \underbrace{\lambda_s^{\,d_G(v_j,v_i)-1}}_{\text{spatial decay}} \cdot \underbrace{\lambda_t^{\,t-\tau-1}}_{\text{temporal decay}} \cdot \underbrace{\cos\big(\mathrm{Enc}(s_k), \mathrm{Enc}(q^{(t)})\big)}_{\text{semantic relevance}},
\end{equation}
trong đó $d_G$ là khoảng cách đường đi ngắn nhất trên đồ thị (lấy cảm hứng từ lý thuyết mạng tổ chức: ảnh hưởng suy giảm theo khoảng cách) và temporal decay phản ánh hiệu ứng gần đây. Các câu có điểm $\ge \theta$ (cùng với câu hỏi hiện tại) được \emph{nhấn mạnh} bằng kỹ thuật attention steering SPA; toàn bộ lịch sử vẫn nằm trong prompt. Siêu tham số: $\lambda_s = \lambda_t = 0.92$, $\theta = 0.65$.

\paragraph{Thực nghiệm.} Năm bộ dữ liệu: HotpotQA, 2WikiMultihopQA, MuSiQue (F1), MATH-500, MMLU-Pro (độ chính xác). Mô hình chính Qwen3-4B; kiểm tra thêm Llama-3.1-8B, Qwen3-32B. Cấu hình 5 agent, 2 vòng (3 vòng với debate), các cấu trúc ngẫu nhiên, phân lớp, đầy đủ và debate.

\begin{table}[!htbp]
\centering
\small
\caption{Kết quả AGENT-RADAR với Qwen3-4B, cấu trúc ngẫu nhiên~\cite{p8radar}.}
\label{tab:radar}
\begin{tabular}{lccccc}
\toprule
\textbf{Phương pháp} & \textbf{HotpotQA} & \textbf{2Wiki} & \textbf{MuSiQue} & \textbf{MATH-500} & \textbf{MMLU-Pro} \\
\midrule
Vanilla (1 agent) & 75.07 & 71.14 & 35.47 & 80.60 & 60.20 \\
CoT-SC ($n=8$)     & 76.51 & 72.03 & 36.33 & 79.80 & 62.60 \\
Compaction         & 74.33 & 73.34 & 34.43 & 83.60 & 56.40 \\
AgentPrune         & 75.31 & 73.03 & 37.05 & 80.80 & 63.60 \\
AgentDropout       & 73.14 & 78.45 & 37.06 & 82.60 & 67.40 \\
\textbf{AGENT-RADAR} & \textbf{80.78} & \textbf{80.81} & \textbf{39.72} & \textbf{88.80} & \textbf{69.40} \\
\bottomrule
\end{tabular}
\end{table}

Khi cắm vào các khung có sẵn, AGENT-RADAR cải thiện GPTSwarm (MuSiQue tăng 12.87 F1), AutoGen (MATH-500 từ 81.0 lên 90.0) và Multi-Agent Debate. Thời gian chạy 1\,196.5 giây/tác vụ so với 968.8 giây của \MAS{} gốc (tăng khoảng 23\%), thấp hơn Compaction (1\,757.1) và AgentDropout (1\,382.4). Thí nghiệm loại bỏ cho thấy chỉ nối các câu được chọn vào prompt (không attention steering) giảm mạnh hiệu năng (MATH-500 từ 88.8 xuống 81.4).

\begin{nhanxet}
\textbf{Điểm mạnh:} không cần huấn luyện, không mất thông tin; kết hợp cấu trúc đồ thị, thời gian và ngữ nghĩa một cách có nguyên lý; cải thiện nhất quán trên 3 khung, 4 cấu trúc và 3 backbone. \textbf{Hạn chế:} cần truy cập logit/chú ý nên không áp dụng được cho mô hình đóng qua API; siêu tham số được điều chỉnh trên MuSiQue -- đồng thời là tập kiểm thử; không có thanh sai số hay kiểm định thống kê dù nhiệt độ sinh bằng 1; không giảm số token; giả định khoảng cách đồ thị phản ánh mức liên quan có thể sai trong hệ có coordinator.
\end{nhanxet}


% ---------------------------------------------------------------------
\subsection{So sánh trong nhóm: DAIN và AGENT-RADAR}
\label{sec:cmp-comm}

Hai bài cùng trả lời câu hỏi ``agent nên nhận thông tin gì từ agent khác'', nhưng ở hai mức trừu tượng rất khác nhau (Bảng~\ref{tab:cmp-g3}).

\begin{table}[!htbp]
\centering
\small
\caption{So sánh DAIN và AGENT-RADAR.}
\label{tab:cmp-g3}
\begin{tabularx}{\textwidth}{L{3.0cm} Y Y}
\toprule
\textbf{Tiêu chí} & \textbf{DAIN} & \textbf{AGENT-RADAR} \\
\midrule
Bản chất agent & Mô-đun nơ-ron (không phải LLM) & LLM agent theo vai trò \\
Kênh giao tiếp & Vector qua low-rank bottleneck & Văn bản (toàn bộ lịch sử hội thoại) \\
Cơ chế chọn lọc & Học được: sparse gate + soft graph $\Phi$ & Tính toán: spatial $\times$ temporal decay $\times$ cosine \\
Cần huấn luyện & Có (đầu-cuối) & Không (plug-in) \\
Đối tượng được chọn & Agent nào hoạt động, ai nói với ai, mạnh bao nhiêu & Câu nào trong lịch sử được nhấn mạnh \\
Hiệu quả tính toán & Giảm (chỉ 3.2--5.5/8 agent chạy) & Tăng nhẹ (+23\% latency), không giảm token \\
Bài toán & Phân loại/hồi quy đa phương thức & QA nhiều bước, toán, MMLU-Pro \\
Kết quả nổi bật & +1.3 đến +2.6 điểm so với MoE & +2.0 đến +7.6 điểm so với AgentDropout \\
\bottomrule
\end{tabularx}
\end{table}

Dù khác nhau về mức trừu tượng, hai bài cùng đi theo một nguyên lý: \emph{lọc mềm} thay vì \emph{cắt cứng}. DAIN gán cường độ liên tục $\Phi_{ij}$ cho từng cặp agent thay vì bật/tắt cố định; AGENT-RADAR nhấn mạnh câu liên quan thay vì xoá câu không liên quan như AgentPrune/AgentDropout. Cả hai cho thấy giữ lại thông tin nhưng điều chỉnh trọng số tốt hơn là loại bỏ hẳn. Khác biệt quan trọng là DAIN đánh đổi tính tổng quát (phải huấn luyện cho từng bài toán) để lấy hiệu quả, còn AGENT-RADAR đánh đổi chi phí và yêu cầu hộp trắng để lấy tính cắm-là-chạy.

% =====================================================================
\section{Nhóm Domain Applications}
% =====================================================================

\subsection{MAC-RRG: Multi-agent collaboration radiology report generation}
\label{sec:macrrg}

\noindent\textit{Wang, Qiao, Wang et al. (ĐH An Huy và cộng sự), arXiv 2609.26124, 09/2026, định dạng IEEE Transactions~\cite{p2macrrg}.}

\paragraph{Bài toán và động lực.} Sinh báo cáo X-quang ngực tự động giúp giảm tải cho bác sĩ chẩn đoán hình ảnh. Các phương pháp chỉ dùng LLM (như R2GenGPT) thiếu tri thức y khoa có cấu trúc nên dễ hallucination; các phương pháp tăng cường bằng knowledge graph (KG) thì kết hợp tri thức \emph{tĩnh, một lần}, không truy vấn lại dựa trên nội dung đã sinh và không dùng tri thức phi cấu trúc. Cả hai đều sinh một lượt, không có vòng phản hồi để kiểm tra thiếu sót.

\paragraph{Phương pháp.} Hai giai đoạn:
\begin{enumerate}
  \item \textbf{Draft generation:} Swin Transformer mã hóa ảnh, Llama2-7B sinh bản nháp $R_0$.
  \item \textbf{Knowledge-guided refinement:} trích thực thể giải phẫu và bệnh lý từ bản nháp bằng từ điển; \emph{MM-KG agent} truy xuất 10 láng giềng hàng đầu của mỗi thực thể trong knowledge graph y khoa (mã hóa bằng Bio\_ClinicalBERT, tổng hợp bằng chú ý đồ thị có nhận thức quan hệ); \emph{Knowledge agent} retrieval đoạn văn từ kho 393 đoạn xây từ PubMedVision (BGE-M3 + reranker). Hai luồng tri thức được chiếu và \emph{ghép nối} với token thị giác làm tiền tố cho LLM sinh báo cáo cuối: $Z = [V; \tilde R; \tilde G]$.
\end{enumerate}

\paragraph{Thực nghiệm và kết quả.} Ba bộ dữ liệu: IU X-Ray, MIMIC-CXR, CheXpert Plus; độ đo BLEU, ROUGE-L, METEOR, CIDEr và clinical efficacy (CE). BLEU-4 đạt 0.197 (IU X-Ray), 0.128 (MIMIC-CXR), 0.105 (CheXpert Plus) -- cao nhất trong các đối chứng được liệt kê. Trên CheXpert Plus, F1 lâm sàng đạt 0.295 so với 0.260 của R2GenGPT. Thí nghiệm loại bỏ cho thấy chỉ cần thêm \emph{một} nguồn tri thức bất kỳ đã nâng F1 từ 0.260 lên khoảng 0.286--0.288; phần cải thiện BLEU-4 nhỏ (0.101 lên 0.105). Trên MIMIC-CXR, phương pháp không dẫn đầu ở ROUGE-L, METEOR và CIDEr.

\begin{nhanxet}
\textbf{Điểm mạnh:} kết hợp tri thức có cấu trúc và phi cấu trúc, neo vào thực thể từ bản nháp có điều kiện hình ảnh; đưa tri thức vào dưới dạng token liên tục thay vì prompt dài; cải thiện clinical efficacy nhất quán; nhiều thí nghiệm siêu tham số. \textbf{Hạn chế:} ``multi-agent'' chủ yếu là cách đặt tên -- các ``agent'' là mô-đun retrieval tất định, không giao tiếp hay thương lượng; ``lặp'' thực tế chỉ một lượt tinh chỉnh; cải thiện NLG nhỏ và không có kiểm định thống kê; checkpoint được chọn dựa trên độ đo đánh giá; CE chỉ báo cáo trên một tập; không có đánh giá của bác sĩ; không báo cáo chi phí.
\end{nhanxet}

% ---------------------------------------------------------------------
\subsection{UrbanAgent: Collaborative reasoning với bằng chứng từ công cụ cho urban region profiling}
\label{sec:urban}

\noindent\textit{Hao, Jiang, Liu et al. (HKUST Quảng Châu và cộng sự), arXiv 2607.13558, KDD 2026~\cite{p6urban}.}

\paragraph{Bài toán và động lực.} Dự đoán chỉ số kinh tế -- xã hội (phát thải carbon, dân số, GDP) của vùng đô thị từ dữ liệu đa phương thức: ảnh vệ tinh, điểm quan tâm (POI), mô tả văn bản, mô hình tòa nhà 3D. Các phương pháp học sâu truyền thống trộn các phương thức vào một biểu diễn ẩn: (1) mâu thuẫn giữa các phương thức (vốn phản ánh sử dụng đất hỗn hợp, không phải nhiễu) bị hấp thụ thay vì được làm rõ; (2) suy luận dựa trên tương quan bề mặt, kém tin cậy ở vùng chưa thấy.

\paragraph{Phương pháp.} Hai giai đoạn:
\begin{enumerate}
  \item \textbf{Tool-augmented agentic learning:} mỗi phương thức một agent (Qwen3-VL-4B) được dạy gọi 4 công cụ: tìm kiếm web, cắt-phóng to ảnh, tra cứu ánh sáng ban đêm, ảnh vệ tinh lịch sử. Huấn luyện SFT khởi động rồi GRPO với phần thưởng tổng hợp
  \begin{equation}
    R_{total} = \alpha R_{outcome} + \beta R_{format} + \gamma R_{process},
  \end{equation}
  trong đó $R_{process}$ phạt việc không dùng công cụ, gọi thừa hoặc gọi sai. Mỗi agent trả về dự đoán và độ tin cậy $c_k$.
  \item \textbf{Suy luận multi-agent collaboration:} xây đồ thị quan hệ dị thể với ba loại cạnh: \emph{corroboration} (dự đoán gần nhau), \emph{rectification} (dự đoán chênh lệch lớn), \emph{supervisory} (agent tin cậy hơn hướng dẫn agent kém tin cậy). Truyền thông điệp có nhận thức quan hệ qua $L=5$ vòng (GRU), sau đó một Hierarchical Spatio-Temporal Final Judge tổng hợp ra dự đoán cuối.
\end{enumerate}

\paragraph{Thực nghiệm và kết quả.} 2\,000 vị trí 1\,km$^2$ trên toàn cầu; nhãn được chuẩn hóa theo hạng; độ đo $R^2$ và Spearman $\rho$. Đối chứng gồm MLLM mã nguồn mở và đóng (GPT-4o, Gemini-2.5-Flash), mô hình tiền huấn luyện đô thị (UrbanCLIP, UrbanVLP) và nhiều chiến lược multi-agent aggregation (AgentFusion, AgentBlender, Agent-Arbiter, Hierarchical Prompting). UrbanAgent đạt $R^2$ khoảng 0.63--0.65 và Spearman khoảng 0.80 trên ba chỉ số trong miền, và $R^2$ khoảng 0.60--0.62 trên các thành phố chưa thấy (Nhật Bản, Na Uy, Madagascar, Mexico). Tác giả báo cáo mức tăng $R^2$ trung bình 8.1\% so với đối chứng mạnh nhất. Bỏ giai đoạn multi-agent gây suy giảm lớn nhất; bỏ công cụ giảm khoảng 13\%; chỉ SFT (không GRPO) cũng kém hơn.

\begin{nhanxet}
\textbf{Điểm mạnh:} mô hình hóa tường minh mâu thuẫn giữa các phương thức bằng quan hệ có kiểu; phần thưởng RL nhắm vào chất lượng quá trình gọi công cụ; chuyển giao tốt sang thành phố mới; so sánh với nhiều chiến lược multi-agent aggregation cho thấy thiết kế cộng tác quan trọng. \textbf{Hạn chế:} giai đoạn ``debate'' là bộ tổng hợp GNN/GRU cần nhãn, không phải debate bằng ngôn ngữ; dữ liệu nhỏ (khoảng 400 điểm kiểm thử) và không có kiểm định thống kê; mô tả văn bản được GPT-4o sinh từ chính ảnh vệ tinh nên các phương thức không độc lập; tìm kiếm web có thể rò rỉ thông tin tiên nghiệm; không có phân tích chi phí.
\end{nhanxet}

% ---------------------------------------------------------------------
\subsection{So sánh trong nhóm: MAC-RRG và UrbanAgent}
\label{sec:cmp-app}

Hai bài ứng dụng ở hai miền khác nhau nhưng có cùng cấu trúc ``nhiều nguồn bằng chứng $\to$ nhiều agent $\to$ tổng hợp''. Bảng~\ref{tab:cmp-g4} cho thấy mức độ ``đa agent'' của hai bài khác nhau đáng kể.

UrbanAgent gần với một \MAS{} thực sự hơn: các agent tự chủ dùng công cụ và mâu thuẫn giữa chúng được mô hình hoá thành quan hệ. MAC-RRG thực chất là một hệ RAG hai nguồn được đặt tên theo khung đa agent. Điều này gợi ý rằng giá trị của cộng tác đa agent trong ứng dụng chuyên ngành đến chủ yếu từ việc \emph{làm rõ và xử lý mâu thuẫn giữa các nguồn bằng chứng}, chứ không chỉ từ việc có nhiều nguồn.

\begin{table}[H]
\centering
\small
\caption{So sánh MAC-RRG và UrbanAgent.}
\label{tab:cmp-g4}
\begin{tabularx}{\textwidth}{L{3.0cm} Y Y}
\toprule
\textbf{Tiêu chí} & \textbf{MAC-RRG} & \textbf{UrbanAgent} \\
\midrule
Miền ứng dụng & Y tế: báo cáo X-quang ngực & Đô thị: carbon, dân số, GDP \\
Dữ liệu đầu vào & Ảnh X-quang + KG + kho văn bản & Ảnh vệ tinh, văn bản, POI, mô hình 3D \\
Bản chất agent & Mô-đun retrieval tất định & VLM dùng công cụ (Qwen3-VL-4B) \\
Agent có tự quyết định? & Không & Có (chọn công cụ, số lượt gọi) \\
Cơ chế tổng hợp & Ghép token tri thức vào tiền tố LLM & Đồ thị quan hệ 3 loại cạnh + GRU + final judge \\
Xử lý mâu thuẫn giữa nguồn & Không có & Tường minh (rectification edge) \\
Huấn luyện & LLM sinh báo cáo & SFT + GRPO cho agent; GNN cho tổng hợp \\
Kiểm tra tổng quát hoá & Không & Có (thành phố chưa thấy) \\
Mức cải thiện & Nhỏ (BLEU-4 +0.004 so với base) & Trung bình (+8.1\% $R^2$) \\
\bottomrule
\end{tabularx}
\end{table}

\chapter{SO SÁNH VÀ ĐÁNH GIÁ CÓ HỆ THỐNG}
\label{chap:comparison}

Nếu Chương~\ref{chap:papers} so sánh từng cặp nghiên cứu trong cùng một nhóm, thì chương này đặt cả tám nghiên cứu cạnh nhau. Các chiều so sánh gồm: phương pháp (\S\ref{sec:cmp-method}), dữ liệu và độ đo (\S\ref{sec:cmp-data}), thiết lập thực nghiệm (\S\ref{sec:cmp-setup}), kết quả và chi phí (\S\ref{sec:cmp-result}), ưu điểm -- hạn chế -- phạm vi áp dụng (\S\ref{sec:cmp-proscons}) và độ tin cậy của bằng chứng (\S\ref{sec:cmp-evidence}). Cuối chương thảo luận các xu hướng xuyên suốt (\S\ref{sec:cmp-trends}).

\section{So sánh về phương pháp}
\label{sec:cmp-method}

\begin{table}[!htbp]
\centering
\footnotesize
\caption{So sánh phương pháp.}
\label{tab:cmp-method}
\begin{tabularx}{\textwidth}{L{1.9cm} L{2.0cm} L{2.4cm} C{1.5cm} C{1.4cm} Y}
\toprule
\textbf{Nghiên cứu} & \textbf{Loại đóng góp} & \textbf{Cơ chế cốt lõi} & \textbf{Cần huấn luyện} & \textbf{Cần hộp trắng} & \textbf{Backbone} \\
\midrule
\#1 MIRAGE & Phòng thủ & Đồ thị ảo + diễn đạt lại & Không & Không & Không nêu trong thân bài \\
\#4 TIDE & Khung đánh giá & Môi trường hidden profile có lừa dối & Không & Không & GPT-5.5, DeepSeek-v4-pro, Grok-4.5 \\
\#3 ProgRouter & Bộ điều phối & Chấm tiến độ + dự báo + Lyapunov & Trực tuyến (nhẹ) & Không & Qwen, Granite, Gemma (0.5B--35B); điều phối 80B \\
\#5 HalluProp & Risk predictor & Noisy-OR khả vi trên đồ thị & Có (giám sát) & Không & GPT-3.5/4o, Llama-3-8B, Qwen-2-7B \\
\#7 DAIN & Kiến trúc kết hợp & Cổng thưa + soft communication graph & Có (đầu-cuối) & --- & ResNet-50, BERT-base \\
\#8 AGENT-RADAR & Mô-đun cắm thêm & Chấm điểm không gian $\times$ thời gian $\times$ ngữ nghĩa + SPA & Không & \textbf{Có} & Qwen3-4B/32B, Llama-3.1-8B \\
\#2 MAC-RRG & Hệ thống ứng dụng & Nháp + 2 nguồn tri thức + tinh chỉnh & Có & --- & Swin, Llama2-7B \\
\#6 UrbanAgent & Hệ thống ứng dụng & Agent dùng công cụ (SFT+GRPO) + GNN quan hệ & Có & --- & Qwen3-VL-4B \\
\bottomrule
\end{tabularx}
\end{table}

\paragraph{Nhận xét.}
\begin{itemize}
  \item \textbf{Không cần huấn luyện vs. cần huấn luyện.} MIRAGE, TIDE và AGENT-RADAR hoạt động không cần huấn luyện, dễ cắm vào hệ có sẵn. ProgRouter chỉ học các cây hồi quy nhẹ theo kiểu trực tuyến. HalluProp, DAIN, MAC-RRG và UrbanAgent cần dữ liệu có nhãn, nên khả năng chuyển giao là câu hỏi mở.
  \item \textbf{Hộp đen vs. hộp trắng.} AGENT-RADAR là phương pháp duy nhất \emph{bắt buộc} truy cập logit/chú ý, nên không áp dụng được với các LLM thương mại qua API -- trong khi phần lớn \MAS{} thực tế dùng loại mô hình này. Ngược lại, MIRAGE, ProgRouter, HalluProp tương thích với mô hình hộp đen.
  \item \textbf{Đối tượng được tối ưu.} ProgRouter tối ưu \emph{việc chọn mô hình}; DAIN tối ưu \emph{việc chọn agent và cường độ kết nối}; AGENT-RADAR tối ưu \emph{việc chọn nội dung để chú ý}; HalluProp chỉ ra \emph{agent/cạnh cần can thiệp}; MIRAGE thay đổi \emph{cái được phơi bày} chứ không thay đổi cái được thực thi.
  \item \textbf{Kỹ thuật lặp lại.} Hai ý tưởng xuất hiện nhiều lần: (a) \emph{truyền thông điệp trên đồ thị} (HalluProp -- Noisy-OR; DAIN -- low-rank bottleneck; UrbanAgent -- GNN quan hệ); (b) \emph{bộ mã hóa câu nhỏ} (MiniLM, BERT) để đo semantic relevance với chi phí thấp (ProgRouter, HalluProp, AGENT-RADAR).
\end{itemize}

\section{So sánh về dữ liệu và thiết lập đánh giá}
\label{sec:cmp-data}

\begin{table}[!htbp]
\centering
\footnotesize
\caption{So sánh dữ liệu, độ đo và độ nghiêm ngặt thực nghiệm.}
\label{tab:cmp-data}
\begin{tabularx}{\textwidth}{L{1.9cm} Y L{2.6cm} C{1.3cm} C{1.9cm}}
\toprule
\textbf{Nghiên cứu} & \textbf{Dữ liệu / benchmark} & \textbf{Độ đo chính} & \textbf{Quy mô MAS} & \textbf{Phương sai / kiểm định} \\
\midrule
\#1 MIRAGE & MMLU, GSM8K, SVAMP, HumanEval (100 tác vụ/tập) & AUC/ACC/F1 của cạnh suy ra & 3--7 agent & Không \\
\#4 TIDE & 120 môi trường Hi-ToM $\times$ 2 điều kiện $\times$ 3 mô hình & Đa số đúng/decoy, evidence lineage, propagation order & 5 (+3) & \textbf{Có} (bootstrap, McNemar) \\
\#3 ProgRouter & HumanEval+, MBPP, MATH-500, ASQA & Độ chính xác, năng lượng J, thời gian & Điều phối + nhiều vai trò & Không \\
\#5 HalluProp & MMLU, MATH, MedQA (500 câu/tập) & AUROC, Hit@k, độ chính xác sau can thiệp & 3--7 & Không \\
\#7 DAIN & ADNI, MIMIC-IV, MM-IMDB, MOSI, ENRICO & Accuracy, MSE, số agent kích hoạt & 8 mô-đun & Nêu có (5 seed) nhưng bảng không có \\
\#8 AGENT-RADAR & HotpotQA, 2Wiki, MuSiQue, MATH-500, MMLU-Pro & F1, accuracy, thời gian & 5--6 & Không \\
\#2 MAC-RRG & IU X-Ray, MIMIC-CXR, CheXpert Plus & BLEU, ROUGE-L, METEOR, CIDEr, CE & 2 mô-đun & Không \\
\#6 UrbanAgent & 2\,000 vị trí toàn cầu + 200 vị trí thành phố mới & $R^2$, Spearman & 4 & Không \\
\bottomrule
\end{tabularx}
\end{table}

\paragraph{Nhận xét.}
\begin{itemize}
  \item \textbf{Benchmark lặp lại:} MMLU, MATH/GSM8K và các bộ QA nhiều bước là ``mặt bằng chung'' của các bài phương pháp (\#1, \#5, \#8). Tuy vậy, mỗi bài dùng một tập con, một cấu trúc và backbone khác nhau nên \emph{không thể so sánh trực tiếp con số giữa các bài}.
  \item \textbf{Quy mô nhỏ:} hầu hết các \MAS{} được thử nghiệm có 3--8 agent và 2--3 vòng. Không bài nào thử nghiệm hệ hàng chục agent hay tác vụ dài hạn (web, lập trình nhiều tệp).
  \item \textbf{Độ nghiêm ngặt thống kê:} chỉ TIDE có kiểm định thống kê đầy đủ. Các bài còn lại không báo cáo phương sai, trong khi một số chạy với nhiệt độ sinh khác 0 (AGENT-RADAR: 1.0; HalluProp: 0.7) và mức cải thiện chỉ 1--3 điểm -- có thể nằm trong khoảng nhiễu.
  \item \textbf{Rủi ro rò rỉ đánh giá:} AGENT-RADAR điều chỉnh siêu tham số trên MuSiQue -- cũng là tập kiểm thử; MAC-RRG chọn checkpoint dựa trên độ đo đánh giá; ProgRouter chỉ báo cáo sau giai đoạn ổn định.
\end{itemize}

\section{So sánh thiết lập thực nghiệm}
\label{sec:cmp-setup}

Để biết các kết quả có so sánh được với nhau hay không, cần đặt cạnh nhau \emph{điều kiện} mà mỗi bài thử nghiệm: dùng model nào, bao nhiêu agent, bao nhiêu vòng, topology ra sao và quy mô dữ liệu đến đâu (Bảng~\ref{tab:cmp-setup}).

\begin{table}[!htbp]
\centering
\footnotesize
\caption{So sánh thiết lập thực nghiệm của tám nghiên cứu.}
\label{tab:cmp-setup}
\begin{tabularx}{\textwidth}{L{1.9cm} Y C{1.75cm} C{1.4cm} L{2.4cm} L{2.2cm}}
\toprule
\textbf{Nghiên cứu} & \textbf{Backbone} & \textbf{Số agent} & \textbf{Số vòng} & \textbf{Topology} & \textbf{Quy mô dữ liệu} \\
\midrule
\#1 MIRAGE & Không nêu trong thân bài & 3--7 & 1 lượt DAG & DAG từ G-Designer, AGP, ARG-Designer & 4 tập $\times$ 100 tác vụ \\
\#4 TIDE & GPT-5.5, DeepSeek-v4-pro, Grok-4.5 & 5 (+3 observer) & T0, T1, R1--R3 & Kết nối đầy đủ + kênh riêng & 120 môi trường $\times$ 2 điều kiện $\times$ 3 model \\
\#3 ProgRouter & Qwen, Granite, Gemma (0.5B--35B); coordinator 80B & 1 coordinator + các worker & Nhiều bước & Phân cấp (ledger) & 164 / 200 / 200 / 100 tác vụ \\
\#5 HalluProp & GPT-3.5, GPT-4o, Llama-3-8B, Qwen-2-7B & 3--7 & 3 & Ngẫu nhiên; sao, chuỗi, phân cấp & 3 tập $\times$ 500 câu \\
\#7 DAIN & ResNet-50, BERT-base (không LLM) & 8 mô-đun & 1 & Soft graph học được & 5 tập dữ liệu \\
\#8 AGENT-RADAR & Qwen3-4B; Llama-3.1-8B, Qwen3-32B & 5 (tối đa 6) & 2--3 & Ngẫu nhiên, phân lớp, đầy đủ, debate & 5 benchmark \\
\#2 MAC-RRG & Swin + Llama2-7B & 2 mô-đun & 1 lượt tinh chỉnh & Dây chuyền cố định & 3 tập X-quang \\
\#6 UrbanAgent & Qwen3-VL-4B & 4 & 5 (GRU) & Đồ thị quan hệ 3 loại cạnh & 2\,000 + 200 vị trí \\
\bottomrule
\end{tabularx}
\end{table}

\paragraph{Nhận xét.}
\begin{itemize}
  \item \textbf{Mô hình nhỏ chiếm ưu thế.} Trừ TIDE dùng các model thương mại lớn, các bài phương pháp chủ yếu thử với model mở cỡ 4B--8B. Điều này hợp lý về chi phí nhưng chưa cho biết phương pháp còn giá trị khi backbone mạnh hơn hay không; AGENT-RADAR là bài hiếm hoi thử thêm model 32B, và mức cải thiện ở đó nhỏ hơn rõ rệt (MuSiQue +3.5 F1 so với +7.1 F1 trên Llama-3.1-8B).
  \item \textbf{Hệ thống nhỏ, hội thoại ngắn.} Không bài nào vượt quá 8 agent hay 5 vòng. Các hiện tượng như lan truyền lỗi (TIDE, HalluProp) hay quá tải ngữ cảnh (AGENT-RADAR) được kỳ vọng mạnh hơn nhiều ở quy mô lớn, nên kết quả hiện tại có thể đang \emph{đánh giá thấp} vấn đề.
  \item \textbf{Topology đa dạng nhưng không thống nhất.} Mỗi bài dùng một kiểu topology khác nhau, nên một phương pháp tốt trên đồ thị ngẫu nhiên (HalluProp, AGENT-RADAR) chưa chắc giữ được ưu thế trên topology phân cấp có điều phối viên như ProgRouter.
\end{itemize}

\section{So sánh về kết quả và chi phí}
\label{sec:cmp-result}

\begin{table}[!htbp]
\centering
\footnotesize
\caption{Kết quả nổi bật và mức độ báo cáo chi phí.}
\label{tab:cmp-result}
\begin{tabularx}{\textwidth}{L{1.9cm} Y L{3.6cm}}
\toprule
\textbf{Nghiên cứu} & \textbf{Kết quả nổi bật} & \textbf{Chi phí được báo cáo} \\
\midrule
\#1 MIRAGE & AUC tấn công giảm xuống 0.64--0.68 (G-Designer), so với 0.70--0.81 của phòng thủ dựa trên prompt & Không \\
\#4 TIDE & Đa số đúng giảm 72.50\% $\to$ 14.17\% khi có một deceiver & Không \\
\#3 ProgRouter & 93.0\% HumanEval+, 79.4\% MBPP trong ngân sách; thua CASCADIA 3.5 điểm trên MATH-500 & \textbf{Có}: năng lượng (J) và thời gian, đo thực \\
\#5 HalluProp & AUROC 84.6; độ chính xác sau can thiệp 67.6\% so với 63.5\% (GUARDIAN) & Có: $<$1 giây/chẩn đoán (không tính huấn luyện) \\
\#7 DAIN & Vượt MoE/MMoE 1.3--2.6 điểm; 3.2--5.5/8 agent kích hoạt & Một phần: số tham số (có mâu thuẫn) \\
\#8 AGENT-RADAR & +7.64 F1 HotpotQA, +6.2 MATH-500 so với AgentDropout; tối đa +12.87 F1 khi cắm vào GPTSwarm & Có: +23\% thời gian; không báo cáo token \\
\#2 MAC-RRG & BLEU-4 cao nhất trên 3 tập; F1 lâm sàng 0.295 so với 0.260 & Không \\
\#6 UrbanAgent & $R^2$ tăng trung bình 8.1\%; chuyển giao tốt sang thành phố mới & Không (chỉ giới hạn 4 lượt gọi công cụ) \\
\bottomrule
\end{tabularx}
\end{table}

\paragraph{Nhận xét.} Chi phí -- một trong những lý do chính khiến \MAS{} khó triển khai -- chỉ được đo nghiêm túc ở ProgRouter (năng lượng GPU thực) và một phần ở AGENT-RADAR (thời gian). Các phương pháp phòng thủ (MIRAGE) và ứng dụng (MAC-RRG, UrbanAgent) đều thêm nhiều lần gọi mô hình hoặc công cụ nhưng không báo cáo chi phí phát sinh. Một quan sát đáng chú ý từ ProgRouter là chiến lược hiệu quả thường dùng mô hình rất nhỏ cho phần lớn các bước (84.3\% lần gọi dùng mô hình 0.5B trên HumanEval+), gợi ý rằng phần lớn công việc trong quy trình multi-agent không đòi hỏi mô hình lớn.

\section{Ưu điểm, hạn chế và phạm vi áp dụng}
\label{sec:cmp-proscons}

\begin{table}[!htbp]
\centering
\footnotesize
\caption{Tổng hợp ưu điểm, hạn chế chính và phạm vi áp dụng.}
\label{tab:cmp-proscons}
\begin{tabularx}{\textwidth}{L{2.0cm} Y Y L{2.8cm}}
\toprule
\textbf{Nghiên cứu} & \textbf{Ưu điểm chính} & \textbf{Hạn chế chính} & \textbf{Phạm vi áp dụng} \\
\midrule
\#1 MIRAGE & Bài toán mới (bí mật cấu trúc); không đổi tính toán thật & Chỉ 1 tấn công, không có kẻ tấn công thích nghi; che giấu chưa hoàn toàn & \MAS{} dạng DAG có phơi bày vết lập luận \\
\#4 TIDE & Thiết kế thực nghiệm chặt; đo quá trình, không chỉ kết quả & Một tác vụ hẹp; không đề xuất phòng thủ & Kiểm toán \MAS{} gộp bằng chứng phân tán \\
\#3 ProgRouter & Routing theo bước; đo năng lượng thực; không cần dữ liệu ngoại tuyến & Mốc tiến độ cần định nghĩa thủ công; bỏ qua chi phí khởi động & Quy trình coordinator -- thực thi với model zoo mở \\
\#5 HalluProp & Pre-hoc, cực nhanh, dễ giải thích & Có giám sát; vai trò bị nhiễu nhân tạo; đồ thị nhỏ & Sàng lọc cấu hình vai trò/cấu trúc trước triển khai \\
\#7 DAIN & Sparse activation; diễn giải được theo PID & Không phải LLM agent; số liệu thiếu nhất quán & Phân loại/hồi quy đa phương thức có giám sát \\
\#8 AGENT-RADAR & Không mất thông tin; không cần huấn luyện; cải thiện nhất quán & Cần hộp trắng; không giảm token; rủi ro rò rỉ siêu tham số & \MAS{} văn bản dùng LLM mã nguồn mở \\
\#2 MAC-RRG & Kết hợp tri thức cấu trúc + phi cấu trúc; cải thiện CE & ``Agent'' chỉ là mô-đun retrieval; cải thiện nhỏ & Sinh văn bản y khoa có neo tri thức \\
\#6 UrbanAgent & Mô hình hóa mâu thuẫn tường minh; RL cho dùng công cụ & Dữ liệu nhỏ; tổng hợp bằng GNN cần nhãn & Dự đoán đa góc nhìn có dữ liệu mâu thuẫn \\
\bottomrule
\end{tabularx}
\end{table}

\section{Đánh giá độ tin cậy của bằng chứng thực nghiệm}
\label{sec:cmp-evidence}

Mục này đánh giá \emph{chất lượng bằng chứng} mà mỗi bài đưa ra, thay vì độ lớn của con số. Nhóm dùng năm tiêu chí: (1) baseline có mạnh và phù hợp không; (2) có thí nghiệm loại bỏ (ablation) để tách đóng góp từng thành phần không; (3) có báo cáo phương sai hoặc kiểm định thống kê không; (4) có đo chi phí không; (5) kết quả có được kiểm tra trên nhiều model, miền hoặc dữ liệu mới không. Mỗi tiêu chí được chấm ba mức: \textbf{$\bullet$} đầy đủ, \textbf{$\circ$} một phần, \textbf{--} không có (Bảng~\ref{tab:cmp-evidence}).

\begin{table}[!htbp]
\centering
\small
\caption{Đánh giá độ tin cậy của bằng chứng thực nghiệm ($\bullet$ đầy đủ, $\circ$ một phần, -- không có).}
\label{tab:cmp-evidence}
\begin{tabular}{l C{1.6cm} C{1.6cm} C{1.9cm} C{1.4cm} C{1.9cm}}
\toprule
\textbf{Nghiên cứu} & \textbf{Baseline} & \textbf{Ablation} & \textbf{Phương sai / kiểm định} & \textbf{Chi phí} & \textbf{Tổng quát hoá} \\
\midrule
\#1 MIRAGE      & $\circ$   & $\bullet$ & --        & --        & $\circ$   \\
\#4 TIDE        & $\circ$   & $\bullet$ & $\bullet$ & --        & $\circ$   \\
\#3 ProgRouter  & $\bullet$ & $\circ$   & --        & $\bullet$ & $\circ$   \\
\#5 HalluProp   & $\circ$   & $\bullet$ & --        & $\circ$   & $\circ$   \\
\#7 DAIN        & --        & $\bullet$ & --        & $\circ$   & $\circ$   \\
\#8 AGENT-RADAR & $\bullet$ & $\bullet$ & --        & $\circ$   & $\bullet$ \\
\#2 MAC-RRG     & $\bullet$ & $\bullet$ & --        & --        & $\circ$   \\
\#6 UrbanAgent  & $\bullet$ & $\circ$   & --        & --        & $\bullet$ \\
\bottomrule
\end{tabular}
\end{table}

\paragraph{Giải thích một số đánh giá.}
\begin{itemize}
  \item \textbf{Baseline.} MIRAGE chỉ so với các phòng thủ prompt-injection chung vì chưa có phòng thủ chuyên biệt cho bài toán mới; TIDE so sánh các điều kiện ghép cặp thay vì phương pháp; HalluProp so với các phương pháp hậu nghiệm vốn là phương pháp \emph{giảm thiểu}, không phải \emph{phát hiện}; DAIN chỉ so với MoE và các cách kết hợp cũ.
  \item \textbf{Ablation.} ProgRouter chỉ làm ablation trên HumanEval+; UrbanAgent chủ yếu trình bày ablation dưới dạng hình, ít số liệu cụ thể.
  \item \textbf{Phương sai.} Đây là điểm yếu chung nhất: chỉ TIDE có kiểm định thống kê đầy đủ. DAIN tuyên bố dùng 5 seed nhưng bảng kết quả không có độ lệch chuẩn.
  \item \textbf{Chi phí.} ProgRouter đo năng lượng GPU thực; AGENT-RADAR đo thời gian; HalluProp đo thời gian chẩn đoán nhưng bỏ qua huấn luyện; DAIN chỉ đếm tham số và phép cộng không khớp.
  \item \textbf{Tổng quát hoá.} AGENT-RADAR thử trên 3 framework, 4 topology và 3 backbone; UrbanAgent thử trên thành phố chưa thấy. Các bài còn lại chỉ thử trong miền.
\end{itemize}

Nhìn tổng thể, AGENT-RADAR và ProgRouter có bằng chứng thực nghiệm cân đối nhất trong các bài phương pháp; TIDE có thiết kế thống kê chặt nhất nhưng phạm vi hẹp; DAIN có bằng chứng yếu nhất do baseline cũ và số liệu không nhất quán. Đánh giá này nên được dùng khi trích dẫn các con số trong báo cáo: những kết quả thuộc các bài có nhiều ô ``--'' cần được xem là xu hướng, chưa phải kết luận chắc chắn.

\section{Các xu hướng và nhận định xuyên suốt}
\label{sec:cmp-trends}

\subsection{Communication graph là đối tượng trung tâm}
Dù mục tiêu khác nhau, sáu trên tám nghiên cứu thao tác trực tiếp trên communication graph: MIRAGE \emph{giấu} nó, HalluProp \emph{phân tích} propagation risk trên nó, AGENT-RADAR dùng \emph{khoảng cách} trên nó để chấm điểm, DAIN và UrbanAgent \emph{học} nó, TIDE \emph{truy vết} dòng chảy bằng chứng trên nó. Điều này khẳng định cách nhìn ``\MAS{} là đồ thị'' đã trở thành khung chuẩn của lĩnh vực, và topology vừa là đòn bẩy hiệu năng, vừa là bề mặt tấn công, vừa là kênh lan truyền rủi ro.

\subsection{Lan truyền là rủi ro cốt lõi -- và thảo luận không tự sửa sai}
TIDE và HalluProp, từ hai góc độ khác nhau, cùng chỉ ra rằng lỗi lan theo các communication edge. TIDE còn cho thấy một kết quả đáng lo ngại: thảo luận nhiều vòng \emph{khuếch đại} chiều hướng của bằng chứng ban đầu thay vì sửa sai, và thông tin sai được tiếp nhận nhiều hơn và lan xa hơn thông tin đúng. Kết quả này thách thức giả định ngầm của nhiều công trình về multi-agent debate rằng ``nhiều agent thảo luận thì sẽ đúng hơn''. Nó cũng gợi ý rằng các phương pháp như AGENT-RADAR -- vốn \emph{tăng cường} sự chú ý vào nội dung liên quan -- cần được đánh giá thêm trong điều kiện có deceiver: nếu thông tin sai cũng ``liên quan'' về ngữ nghĩa, liệu nó có bị khuếch đại thêm?

\subsection{Đánh đổi chất lượng -- chi phí}
ProgRouter, DAIN và AGENT-RADAR đều đặt câu hỏi về hiệu quả nhưng theo các cách khác nhau: chọn mô hình rẻ hơn, kích hoạt ít agent hơn, hay chú ý đúng chỗ hơn mà không cắt bớt ngữ cảnh. Một điểm chung là \emph{không cần dùng toàn bộ năng lực ở mọi lúc}: ProgRouter dùng chủ yếu mô hình nhỏ; DAIN chỉ dùng khoảng một nửa số agent; hiệu năng của DAIN đạt đỉnh ở ngân sách $\tau = 4$ chứ không phải $8$; AGENT-RADAR tốt nhất ở 5 agent chứ không phải nhiều hơn.

\subsection{Mâu thuẫn giữa minh bạch và bảo mật}
TIDE kết luận rằng hệ thống cần \emph{xác minh nguồn gốc bằng chứng và đường lan truyền một cách tường minh}, tức cần minh bạch về ai nói gì với ai. Ngược lại, MIRAGE cố tình \emph{che giấu và làm nhiễu} chính thông tin đó với người quan sát bên ngoài. Hai hướng không mâu thuẫn hoàn toàn (một bên là kiểm toán nội bộ, một bên là phơi bày ra ngoài), nhưng cho thấy thiết kế \MAS{} an toàn cần phân biệt rõ các ``góc nhìn'' khác nhau cho các bên khác nhau -- đúng như cơ chế dual-view của MIRAGE.

\subsection{Khái niệm ``agent'' bị nới rộng}
Như Bảng~\ref{tab:agentnature} cho thấy, ba trên tám bài (DAIN, MAC-RRG và một phần UrbanAgent) gọi các mô-đun nơ-ron hoặc mô-đun retrieval là ``agent''. Xu hướng gắn nhãn ``multi-agent'' cho các kiến trúc vốn là MoE hay RAG làm mờ ranh giới của lĩnh vực và gây khó cho việc so sánh. Nhóm đề xuất phân biệt rõ ít nhất ba cấp: (i) \emph{LLM agent tự trị} giao tiếp bằng ngôn ngữ; (ii) \emph{LLM agent có tổng hợp học được}; (iii) \emph{mô-đun chuyên biệt} được gọi là agent theo nghĩa ẩn dụ.

\subsection{Ứng dụng chuyên ngành đi theo hướng ``học'' cộng tác}
Hai bài ứng dụng (MAC-RRG, UrbanAgent) đều thay thế debate bằng ngôn ngữ bằng các cơ chế tổng hợp \emph{được huấn luyện} (ghép token, GNN). Lý do khả dĩ là trong các miền như y tế hay địa lý, đầu ra là số hoặc văn bản có cấu trúc chặt, có dữ liệu nhãn, và yêu cầu độ chính xác cao hơn tính linh hoạt. Điều này gợi ý một khoảng cách giữa ``\MAS{} tổng quát'' (không huấn luyện, giao tiếp bằng ngôn ngữ) và ``\MAS{} chuyên ngành'' (huấn luyện đầu-cuối).

\chapter{KHOẢNG TRỐNG NGHIÊN CỨU VÀ HƯỚNG PHÁT TRIỂN}
\label{chap:gaps}

\section{Các khoảng trống nghiên cứu}

\paragraph{G1. Thiếu chuẩn đánh giá thống nhất cho MAS.} Mỗi nghiên cứu tự chọn tập con dữ liệu, topology, số agent và backbone. Không có benchmark chung đo đồng thời \emph{chất lượng}, \emph{chi phí} và \emph{robustness}, nên không thể đặt các phương pháp như ProgRouter, AGENT-RADAR hay HalluProp cạnh nhau để so sánh.

\paragraph{G2. Độ nghiêm ngặt thống kê thấp.} Bảy trên tám bài không báo cáo phương sai hay kiểm định ý nghĩa, dù \MAS{} vốn có tính ngẫu nhiên cao (nhiều lần lấy mẫu, nhiều agent). Các mức cải thiện 1--3 điểm vì vậy khó được khẳng định.

\paragraph{G3. Quy mô nhỏ và tác vụ ngắn.} Các hệ được thử nghiệm có tối đa khoảng 8 agent và 2--5 vòng, chủ yếu trên trắc nghiệm, toán và QA. Hành vi ở quy mô hàng chục agent, tác vụ dài hạn có công cụ (duyệt web, lập trình nhiều tệp) gần như chưa được khảo sát. Nhiều kết quả có thể không còn đúng khi lịch sử dài hơn hoặc đồ thị lớn hơn.

\paragraph{G4. Thiếu kẻ tấn công thích nghi và thiếu phòng thủ.} MIRAGE chỉ được thử với một tấn công cố định; TIDE đo tác hại nhưng không đề xuất phòng thủ; observer chỉ làm nhóm thận trọng hơn chứ không sửa được lỗi. Chưa có phương pháp nào chứng minh được khả năng \emph{khôi phục sự thật} khi có deceiver.

\paragraph{G5. Tương tác giữa các nhóm chưa được nghiên cứu.} Các phương pháp tối ưu hiệu năng (AGENT-RADAR, ProgRouter) chưa được đánh giá dưới điều kiện đối kháng: cơ chế tăng cường nội dung ``liên quan'' có thể vô tình khuếch đại thông tin sai; router tiết kiệm có thể giao bước then chốt cho mô hình nhỏ dễ bị lừa. Ngược lại, các phương pháp an toàn (MIRAGE) chưa đo chi phí phát sinh.

\paragraph{G6. Phụ thuộc vào thiết kế thủ công và dữ liệu có nhãn.} ProgRouter cần mốc tiến độ cho từng miền; DAIN cần phân bổ số agent theo loại thủ công; UrbanAgent cần ngưỡng cạnh chọn tay; HalluProp cần nhãn từ các lần chạy thật. Khả năng tổng quát hóa sang miền hoặc mô hình mới chưa được kiểm chứng.

\paragraph{G7. Phụ thuộc vào truy cập hộp trắng.} Các kỹ thuật attention steering (AGENT-RADAR) không áp dụng được cho mô hình thương mại -- vốn chiếm phần lớn triển khai thực tế.

\paragraph{G8. Thiếu định nghĩa thống nhất về ``agent''.} Như đã phân tích ở \S\ref{sec:cmp-trends}, việc gọi các mô-đun nơ-ron hay mô-đun retrieval là ``agent'' làm khó việc tích lũy tri thức trong lĩnh vực.

\section{Đề xuất hướng phát triển}

\begin{enumerate}
  \item \textbf{Benchmark tổng hợp ba trục.} Xây dựng bộ đánh giá trong đó mỗi cấu hình \MAS{} được đo đồng thời độ chính xác, chi phí (token/năng lượng/độ trễ) và độ suy giảm khi có agent đối kháng (theo giao thức ghép cặp C0/C1 của TIDE). Đây là hướng khả thi nhất cho quy mô một bài tập lớn: có thể kết hợp các tác vụ của TIDE với các cấu trúc ngẫu nhiên của AGENT-RADAR/HalluProp.

  \item \textbf{Giao tiếp có nhận thức nguồn gốc.} Mở rộng điểm số của AGENT-RADAR bằng một thành phần \emph{độ tin cậy nguồn}: phân biệt bằng chứng trực tiếp với thông tin được chuyển tiếp, giảm trọng số theo propagation order (đúng với phát hiện của TIDE rằng thông tin sai lan qua chuỗi dài hơn). Có thể kết hợp với intrinsic risk $\pi_i$ của HalluProp làm trọng số cho từng agent nguồn:
  \begin{equation}
    \text{score}'(s_k) = \text{score}(s_k)\cdot(1 - \hat h_j)\cdot \mu^{\,\text{hop}(s_k)},\qquad 0<\mu<1.
  \end{equation}

  \item \textbf{Điều phối nhận thức rủi ro.} Kết hợp ProgRouter với HalluProp: hệ số phạt trong hàm điểm routing không chỉ dựa trên chi phí mà còn trên rủi ro hallucination dự báo của vai trò, sao cho các bước có rủi ro cao được giao cho mô hình mạnh hơn hoặc thêm bước kiểm chứng.

  \item \textbf{Phòng thủ chủ động trước thông tin sai.} Thay vì observer thụ động, thiết kế \emph{verifier agent} có nhiệm vụ yêu cầu bằng chứng trực tiếp và kiểm tra tính độc lập của nguồn trước khi đồng thuận; đánh giá trên môi trường TIDE.

  \item \textbf{Chuẩn hóa thuật ngữ và báo cáo.} Đề xuất danh mục kiểm tra khi công bố \MAS: bản chất agent, backbone, số agent và vòng, topology, số lần chạy và phương sai, chi phí. Điều này giúp các nghiên cứu sau so sánh được với nhau.

  \item \textbf{Mở rộng quy mô và tác vụ.} Đánh giá các phương pháp trên hệ lớn hơn (hàng chục agent), tác vụ dài hạn có công cụ, và với cả mô hình thương mại -- ví dụ thay SPA bằng các phương pháp nhấn mạnh ở mức prompt cho mô hình hộp đen.
\end{enumerate}

\chapter{KẾT LUẬN}
\label{chap:conclusion}

Báo cáo đã khảo sát tám nghiên cứu mới nhất (tháng 5 -- tháng 9/2026) về multi-agent system và collaborative reasoning, phân loại chúng thành bốn hướng: safety \& robustness (MIRAGE, TIDE), orchestration \& control (ProgRouter, HalluProp), communication \& interaction (DAIN, AGENT-RADAR), và domain applications (MAC-RRG, UrbanAgent). Với mỗi nghiên cứu, báo cáo đã trình bày bài toán, phương pháp, thực nghiệm, kết quả và đánh giá phản biện; sau đó so sánh có hệ thống theo phương pháp, dữ liệu, kết quả, chi phí, ưu nhược điểm và phạm vi áp dụng.

Các nhận định chính gồm:
\begin{itemize}
  \item \textbf{Communication graph} là đối tượng trung tâm của lĩnh vực: nó vừa quyết định hiệu năng, vừa là kênh lan truyền rủi ro, vừa là tài sản cần bảo vệ.
  \item \textbf{Lan truyền} lỗi và thông tin sai là rủi ro cốt lõi; thảo luận nhiều vòng có xu hướng khuếch đại hơn là sửa sai, đi ngược giả định phổ biến rằng cộng tác luôn làm hệ thống chính xác hơn.
  \item \textbf{Hiệu quả} không đến từ việc dùng nhiều agent hơn hay mô hình lớn hơn, mà từ việc kích hoạt đúng agent, đúng mô hình và chú ý đúng nội dung.
  \item \textbf{Độ nghiêm ngặt thực nghiệm} của lĩnh vực còn hạn chế: quy mô nhỏ, thiếu kiểm định thống kê, thiếu đo chi phí và thiếu định nghĩa thống nhất về ``agent''.
\end{itemize}

Từ đó, báo cáo đề xuất các hướng phát triển: xây dựng benchmark đo đồng thời chất lượng -- chi phí -- robustness, thiết kế cơ chế giao tiếp có nhận thức nguồn gốc bằng chứng, kết hợp điều phối với dự báo rủi ro, phát triển phòng thủ chủ động trước thông tin sai và chuẩn hóa cách báo cáo kết quả \MAS. Nhóm cho rằng giai đoạn tiếp theo của lĩnh vực sẽ chuyển từ câu hỏi ``nhiều agent có tốt hơn một agent không'' sang câu hỏi ``làm thế nào để một multi-agent system vừa hiệu quả, vừa đáng tin cậy và vừa an toàn''.

\chapter{ĐÓNG GÓP CỦA CÁC THÀNH VIÊN}
\label{chap:contrib}

% Bảng phân công dưới đây dựa trên cách chia 2 bài/người đã thống nhất.
% Nhóm điều chỉnh phần "Phần chung" và tỉ lệ đóng góp cho đúng thực tế.

Mỗi thành viên phụ trách đọc, phân tích và viết phần trình bày chi tiết cho hai nghiên cứu cùng hướng, đồng thời đảm nhận một phần chung của báo cáo. Bảng~\ref{tab:contrib} tóm tắt phân công.

\begin{table}[!htbp]
\centering
\small
\caption{Phân công công việc trong nhóm.}
\label{tab:contrib}
\begin{tabularx}{\textwidth}{L{3.0cm} L{3.6cm} Y C{1.4cm}}
\toprule
\textbf{Thành viên} & \textbf{Nghiên cứu phụ trách} & \textbf{Phần chung} & \textbf{Đóng góp} \\
\midrule
Trần Xuân Hương & \#1 MIRAGE, \#4 TIDE (Safety \& Robustness) & Tìm và chọn bài; Chương~1 (Giới thiệu); tổng hợp, biên tập báo cáo & 25\% \\
Vũ Tùng Dương & \#3 ProgRouter, \#5 HalluProp (Orchestration \& Control) & Chương~2 (Kiến thức nền); Chương~6 (Khoảng trống \& hướng phát triển) & 25\% \\
Phạm Việt Dũng & \#7 DAIN, \#8 AGENT-RADAR (Communication \& Interaction) & Chương~3 (Phân loại); Chương~5 (So sánh \& đánh giá) & 25\% \\
Lê Văn Đức & \#2 MAC-RRG, \#6 UrbanAgent (Domain Applications) & Chương~7 (Kết luận); thiết kế slide trình bày & 25\% \\
\bottomrule
\end{tabularx}
\end{table}

Tất cả thành viên cùng tham gia thảo luận khung phân loại, rà soát chéo số liệu giữa các phần và chuẩn bị thuyết trình.


\begin{thebibliography}{99}
\addcontentsline{toc}{chapter}{TÀI LIỆU THAM KHẢO}
\setlength{\itemsep}{2pt}

% ----- Tám nghiên cứu được khảo sát -----
\bibitem{p1mirage}
L.~He, Z.~Wen, S.~Su, X.~Li, X.~Wang, \emph{et al.},
``Concealing LLM-based multi-agent topology via phantom structure injection,''
\emph{arXiv preprint arXiv:2609.37567}, 2026.

\bibitem{p2macrrg}
F.~Wang, Y.~Qiao, X.~Wang, \emph{et al.},
``MAC-RRG: Iterative multi-agent collaboration for X-ray radiology report generation,''
\emph{arXiv preprint arXiv:2609.26124}, 2026.

\bibitem{p3progrouter}
S.~Li, A.~M.~Abdelmoniem, and S.~Wang,
``ProgRouter: Online progress-guided orchestration for multi-agent LLM workflows under quality-cost tradeoffs,''
\emph{arXiv preprint arXiv:2608.25992}, 2026.

\bibitem{p4tide}
C.~Yan, Z.~Yue, F.~Zhao, \emph{et al.},
``When truth is distributed: Misinformation derails collective fact recovery in LLM-based multi-agent systems,''
\emph{arXiv preprint arXiv:2608.03421}, 2026.

\bibitem{p5halluprop}
S.~Lin, C.~Wang, P.~Qian, \emph{et al.},
``Before agents speak: Pre-hoc failure risk inference in multi-agent systems,''
\emph{arXiv preprint arXiv:2607.26836}, 2026.

\bibitem{p6urban}
X.~Hao, Y.~Jiang, J.~Liu, \emph{et al.},
``Multi-agent collaborative reasoning with tool-augmented evidence for urban region profiling,''
in \emph{Proc. 32nd ACM SIGKDD Conf. Knowledge Discovery and Data Mining (KDD)}, 2026. arXiv:2607.13558.

\bibitem{p7dain}
X.~Chen, Y.~Li, Z.~Wang, \emph{et al.},
``DAIN: Dynamic agent-based interaction network for efficient and collaborative multimodal reasoning,''
\emph{arXiv preprint arXiv:2606.30189}, 2026.

\bibitem{p8radar}
H.~Zhang, Y.~Tian, and T.~Zhang,
``Enhancing multi-agent communication through attention steering with context relevance,''
\emph{arXiv preprint arXiv:2605.30136}, 2026.

% ----- Tài liệu nền -----
\bibitem{camel}
G.~Li, H.~A.~A.~K.~Hammoud, H.~Itani, D.~Khizbullin, and B.~Ghanem,
``CAMEL: Communicative agents for `mind' exploration of large language model society,''
in \emph{Advances in Neural Information Processing Systems (NeurIPS)}, 2023.

\bibitem{autogen}
Q.~Wu, G.~Bansal, J.~Zhang, \emph{et al.},
``AutoGen: Enabling next-gen LLM applications via multi-agent conversation,''
\emph{arXiv preprint arXiv:2308.08155}, 2023.

\bibitem{metagpt}
S.~Hong, M.~Zhuge, J.~Chen, \emph{et al.},
``MetaGPT: Meta programming for a multi-agent collaborative framework,''
in \emph{Int. Conf. Learning Representations (ICLR)}, 2024.

\bibitem{du2024debate}
Y.~Du, S.~Li, A.~Torralba, J.~B.~Tenenbaum, and I.~Mordatch,
``Improving factuality and reasoning in language models through multiagent debate,''
in \emph{Int. Conf. Machine Learning (ICML)}, 2024.

\bibitem{gptswarm}
M.~Zhuge, W.~Wang, L.~Kirsch, F.~Faccio, D.~Khizbullin, and J.~Schmidhuber,
``GPTSwarm: Language agents as optimizable graphs,''
in \emph{Int. Conf. Machine Learning (ICML)}, 2024.

\bibitem{gdesigner}
G.~Zhang, Y.~Yue, X.~Sun, \emph{et al.},
``G-Designer: Architecting multi-agent communication topologies via graph neural networks,''
in \emph{Int. Conf. Machine Learning (ICML)}, 2025.

\bibitem{agentprune}
G.~Zhang, Y.~Yue, Z.~Li, \emph{et al.},
``Cut the crap: An economical communication pipeline for LLM-based multi-agent systems,''
in \emph{Int. Conf. Learning Representations (ICLR)}, 2025.

\bibitem{magenticone}
A.~Fourney, G.~Bansal, H.~Mozannar, \emph{et al.},
``Magentic-One: A generalist multi-agent system for solving complex tasks,''
\emph{arXiv preprint arXiv:2411.04468}, 2024.

\bibitem{react}
S.~Yao, J.~Zhao, D.~Yu, N.~Du, I.~Shafran, K.~Narasimhan, and Y.~Cao,
``ReAct: Synergizing reasoning and acting in language models,''
in \emph{Int. Conf. Learning Representations (ICLR)}, 2023.

\bibitem{lostmiddle}
N.~F.~Liu, K.~Lin, J.~Hewitt, A.~Paranjape, M.~Bevilacqua, F.~Petroni, and P.~Liang,
``Lost in the middle: How language models use long contexts,''
\emph{Transactions of the Association for Computational Linguistics}, vol.~12, pp.~157--173, 2024.

\bibitem{stasser1985}
G.~Stasser and W.~Titus,
``Pooling of unshared information in group decision making: Biased information sampling during discussion,''
\emph{Journal of Personality and Social Psychology}, vol.~48, no.~6, pp.~1467--1478, 1985.

\end{thebibliography}


\end{document}
